\documentclass[11pt]{amsart}

\usepackage[utf8]{inputenc}
\usepackage[T1]{fontenc}
\IfFileExists{lmodern.sty}{\usepackage{lmodern}}{}
% Unicode characters occurring inside Lean identifiers
\DeclareUnicodeCharacter{03C0}{\ensuremath{\pi}}
\DeclareUnicodeCharacter{03B9}{\ensuremath{\iota}}
\DeclareUnicodeCharacter{2964}{\ensuremath{\Rightarrow}}
\DeclareUnicodeCharacter{22A5}{\ensuremath{\bot}}
\usepackage{amsmath,amssymb,amsthm,mathtools}
\usepackage[margin=1.15in]{geometry}
\usepackage{booktabs,tabularx,longtable}
\usepackage{xcolor}
\usepackage{listings}
\usepackage[obeyspaces]{url}
\usepackage[colorlinks=true,linkcolor=blue!50!black,citecolor=blue!50!black,urlcolor=blue!50!black]{hyperref}

% Lean identifiers and file names (typewriter, breakable, underscores allowed)
\DeclareUrlCommand\lean{\urlstyle{tt}}
\makeatletter
\g@addto@macro\UrlBreaks{\do\-\do\_\do\.\do\A\do\B\do\C\do\D\do\E\do\F\do\G\do\H\do\I\do\J\do\K\do\L\do\M\do\N\do\O\do\P\do\Q\do\R\do\S\do\T\do\U\do\V\do\W\do\X\do\Y\do\Z}
\makeatother
\Urlmuskip=0mu plus 1mu
\setlength{\emergencystretch}{3em}

% Space between prose paragraphs; no paragraph indentation
\newcommand{\spa}{\medskip}
\setlength{\parindent}{0pt}

% Notation
\newcommand{\Ocris}{\mathcal{O}_{\mathrm{cris}}}
\newcommand{\N}{\mathbb{N}}
\newcommand{\Q}{\mathbb{Q}}
\newcommand{\F}{\mathbb{F}}
\DeclareMathOperator{\Spec}{Spec}
\DeclareMathOperator{\Hom}{Hom}
\DeclareMathOperator{\Map}{Map}
\DeclareMathOperator{\Tot}{Tot}
\DeclareMathOperator{\Coind}{Coind}
\newcommand{\INF}{\mathrm{INF}}
\newcommand{\Mod}{\mathrm{Mod}}
\newcommand{\dpow}[1]{^{[#1]}}
\newcommand{\pow}{\mathrm{pow}}
% Public repository (to be filled in before release)
\newcommand{\repourl}{https://github.com/daddeziomath/affine-crystalline}

\theoremstyle{plain}
\newtheorem{theorem}{Theorem}[section]
\newtheorem{proposition}[theorem]{Proposition}
\newtheorem{lemma}[theorem]{Lemma}
\newtheorem{corollary}[theorem]{Corollary}
\newtheorem*{theoremA}{Theorem A}
\newtheorem*{theoremB}{Theorem B}
\theoremstyle{definition}
\newtheorem{definition}[theorem]{Definition}
\newtheorem{setting}[theorem]{Setting}
\newtheorem{example}[theorem]{Example}
\theoremstyle{remark}
\newtheorem{remark}[theorem]{Remark}

% "Lean:" pointer after a statement
\newcommand{\leanref}[1]{\par\smallskip\noindent{\small\textit{Lean:} #1}\par\smallskip}

\lstdefinelanguage{lean4}{
  morekeywords={theorem,lemma,def,structure,where,fun,by,Type,Nonempty},
  sensitive=true,
  morecomment=[l]{--},
  morestring=[b]",
}
\lstset{
  language=lean4,
  basicstyle=\ttfamily\small,
  keywordstyle=\bfseries\color{blue!60!black},
  commentstyle=\itshape\color{gray},
  columns=fullflexible,
  keepspaces=true,
  breaklines=true,
  mathescape=true,
  frame=single,
  framesep=4pt,
  xleftmargin=6pt,
  xrightmargin=6pt,
}

\title{Formalising crystalline de Rham comparison\\ on a bounded infinitesimal site}
\author{Marco D'Addezio}
\address{Université de Strasbourg}
\email{daddezio@unistra.fr}
\date{October 22, 2026}

\begin{document}

\begin{abstract}
We present a complete Lean~4 formalisation of the affine crystalline de Rham comparison theorem. For this scope, we replace the standard crystalline site with a variant of the infinitesimal site endowed with a crystalline structural sheaf. On the way, we formalise the existence of divided power envelopes, building on the formalisation of divided powers by Chambert-Loir and de Frutos-Fernández. The Lean code was produced with the AI system Claude.
\end{abstract}

\maketitle

\noindent\textbf{Code availability.} The formalisation and this article are available at \expandafter\url\expandafter{\repourl}. The formalisation is archived on Zenodo with DOI \href{https://doi.org/10.5281/zenodo.23573534}{10.5281/zenodo.23573534}.

\setcounter{tocdepth}{1}
\tableofcontents

%======================================================================
\section{Introduction}\label{sec:intro}
%======================================================================

Let $p$ be a prime, let $(A, I, \gamma)$ be a divided power ring with $p \in I$ and $p$ nilpotent in $A$, and let $C$ be an $A$-algebra with $I \cdot C = 0$; write $X \coloneqq \Spec(C)$ and $S \coloneqq \Spec(A)$. If $P \twoheadrightarrow C$ is a surjection from a polynomial $A$-algebra and $D$ is the divided power envelope of its kernel, compatible with $\gamma$, then
\[
R\Gamma\big(\mathrm{Cris}(X/S), \mathcal{O}_{X/S}\big) \;\simeq\; \Omega^\bullet_{D},
\]
where $\Omega^\bullet_D$ is the de Rham complex of $D$ with divided power differentials \cite{berthelot,berthelot-ogus,bhatt-dejong}, \cite[Tag~07LG]{stacks}. This article reports a complete formalisation in Lean~4 with Mathlib \cite{mathlib} of an affine version of this comparison, with trivial coefficients, in which the cohomology is computed on an infinitesimal site instead of the crystalline site.

\spa

The main difference with the classical approach is the site. Instead of the crystalline site, whose objects carry divided powers, we use an \emph{infinitesimal} site whose objects are surjections $B \twoheadrightarrow C$ whose kernel has a uniform nilpotence exponent. Divided powers then enter only through the structure sheaf $\Ocris$, which sends a thickening to its divided power envelope and is used only through the universal property of envelopes. The precise hypotheses are collected in Setting~\ref{Setting:s}, and the Frobenius levels $\INF_e$, whose objects are the thickenings with kernel killed by $x \mapsto x^{p^e}$, the Čech objects $X_e(n)$ and the algebras $D(n) \coloneqq \Ocris(X_e(n))$ are introduced in Sections~\ref{sec:site} and~\ref{sec:ocris}.

\spa

The cohomology $R\Gamma(\INF(C/A), \Ocris)$ in Theorems~A and~B is the cohomology of this infinitesimal site. We prove in Lean that it is already computed on $\INF_e$ for every $e$ with $(p^e)! = 0$ in $A$, for instance for $e = N$ if $p^N = 0$ in $A$ (Theorem~\ref{thm:adjunction-derived}): the truncation functor $\INF(C/A) \to \INF_e$ is left adjoint to the inclusion, and it does not change divided power envelopes. That $R\Gamma(\INF(C/A), \Ocris)$ agrees with the crystalline cohomology $R\Gamma(\mathrm{Cris}(X/S), \mathcal{O}_{X/S})$ above is classical: it follows from the comparison of both with $\Omega^\bullet_D$, or from the description of affine crystalline cohomology through affine divided power thickenings \cite[Tag~07JK]{stacks}. This agreement is not formalised (Section~\ref{sec:scope}), and the library does not define the crystalline site. Berthelot's crystalline site therefore appears in this article only as background and motivation.

\spa

The two main results are the following.

\begin{theoremA}
Assume Setting~\ref{Setting:s}, so that $e$ is a positive integer with $(p^e)! = 0$ in $A$, and let $E$ be a family of compatible divided power envelopes. For every $n \ge 0$ there is an isomorphism of $A$-modules
\[
H^n\big(R\Gamma(\INF(C/A), \Ocris)\big) \;\simeq\; H^n\big(D(0) \to D(1) \to D(2) \to \cdots\big).
\]
\end{theoremA}

\begin{theoremB}
In the situation of Theorem~A, for every $k \ge 0$ there is an isomorphism of $A$-modules
\[
H^k\big(R\Gamma(\INF(C/A), \Ocris)\big) \;\simeq\; H^k\big(\Omega^\bullet_{D(0)}\big),
\]
where $D(0) = \Ocris(X_e(0))$ satisfies the universal property of a compatible divided power envelope of $A[x_c : c \in C]$ with respect to the kernel of $x_c \mapsto c$ (Proposition~\ref{prop:cech-envelope}).
\end{theoremB}

Theorem~A is Theorem~\ref{thm:cech-alexander} and Theorem~B is Theorem~\ref{thm:main}. Both are proved in Lean without \texttt{sorry}, and both use only the axioms \lean{propext}, \lean{Classical.choice} and \lean{Quot.sound}. Both are stated for an arbitrary family $E$, and such a family exists: divided power envelopes exist (Theorem~\ref{thm:envelope-existence}), and this is formalised as well. The theorems are therefore unconditional.

\spa

Section~\ref{sec:motivation} explains the choice of site. Section~\ref{sec:site} develops the site and Section~\ref{sec:ocris} the sheaf $\Ocris$. Section~\ref{sec:pd-derham} constructs divided power differentials and de Rham complexes, Section~\ref{sec:homological} collects the homological algebra, Section~\ref{sec:level} proves Theorem~A, and Section~\ref{sec:derham} proves Theorem~B. Section~\ref{sec:produced} explains how the formalisation was produced, Section~\ref{sec:scope} what remains, and Appendix~\ref{app:index} indexes the main declarations.

\spa

Every statement in this article followed by a \emph{Lean} pointer is proved in the library without \texttt{sorry}. Names are given up to their namespace. Most are in \lean{AffineCrystalline} or one of its sub-namespaces, such as \lean{AffineCrystalline.InfinitesimalSite}; the general material on envelopes, divided power differentials and de Rham complexes, such as \lean{IsDPEnvelope}, \lean{DPEnvelopeChoice}, \lean{PDKaehler} and \lean{DeRhamComplex}, is in the root namespace or in namespaces of Mathlib. Paths are relative to the folder \lean{AffineCrystalline/} of the repository.

%======================================================================
\section{Motivation: an infinitesimal site without divided powers}\label{sec:motivation}
%======================================================================

Berthelot's crystalline site appears in this article only as background and motivation: the cohomology we formalise is defined on the infinitesimal site of Section~\ref{sec:site}, and its relation with Berthelot's site is recalled in Section~\ref{subsec:comparison}.

\subsection{Why the objects carry no divided powers}

In Lean a divided power structure is data indexed by an ideal: its type is \lean{DividedPowers K} for a specific ideal $K$. Rewriting $K$ along an equality of ideals needs transport lemmas, and every map built along the way needs a separate compatibility proof. A Čech computation needs coproducts and Čech objects, and each of them would have to be equipped with divided powers and checked by hand.

\spa

Our idea is to avoid divided power structures on the objects of the site altogether. Instead, we impose a condition on the infinitesimal thickenings themselves, a uniform bound on the nilpotence exponent of the kernel; it is a property, not extra data, and only involves ordinary commutative algebra.

\begin{definition}\label{def:thickening}
An object of the \emph{bounded infinitesimal site} is a surjective $A$-algebra map $\pi \colon B \twoheadrightarrow C$, with no divided powers, whose kernel has a uniform nilpotence exponent: there is an integer $n \ge 0$ such that $x^{n} = 0$ for every $x \in \ker(\pi)$.
A morphism $(B, \pi) \to (B', \pi')$ is an $A$-algebra map $f \colon B \to B'$ with $\pi' \circ f = \pi$, with no further condition.
\end{definition}

The exponent condition is the field \lean{kernel_bounded} of the structure \lean{InfinitesimalThickening}. The use of this site is taken from \cite{edged}.

\subsection{What the exponent condition provides}

It gives truncations, levels, binary coproducts and a weakly initial object in each level, all in \lean{InfinitesimalSite/Filtration.lean} and described in Section~\ref{sec:site}. None of them requires building or checking a divided power structure. Divided powers enter only when $\Ocris$ is applied, and only through the universal property \lean{IsDPEnvelope}.

\subsection{Nothing is lost when \texorpdfstring{$p$}{p} is nilpotent}

The restriction to bounded exponents costs nothing when $p$ is nilpotent.

\begin{lemma}\label{lem:pd-exponent}
Let $n$ be an integer with $n! = 0$ in $A$, let $D$ be an $A$-algebra and let $K \subseteq D$ be an ideal with divided powers $\delta$. Then $x^{n} = 0$ for every $x \in K$.
\end{lemma}

\begin{proof}
By induction on $m$, using $\delta_1(x) = x$ and $\delta_m(x)\delta_k(x) = \binom{m+k}{m}\delta_{m+k}(x)$, one gets $x^m = m!\,\delta_m(x)$. Hence $x^{n} = n!\,\delta_{n}(x)$, which vanishes because $n! = 0$ in $A$.
\end{proof}

\leanref{\lean{dividedPowers_pow_eq_zero_of_factorial_eq_zero} in \lean{InfinitesimalSite/Envelope.lean}; the identity $x^m = m!\,\delta_m(x)$ is Mathlib's \lean{DividedPowers.factorial_mul_dpow_eq_pow}.}

Neither $p$ nor a compatibility of $\delta$ with $\gamma$ is needed. When $p$ is nilpotent, the hypothesis of the lemma holds for suitable $n$.

\begin{corollary}\label{cor:pd-exponent-nilpotent}
Let $N \ge 0$ be an integer with $p^N = 0$ in $A$. Then $(p^N)! = 0$ in $A$. In particular, $x^{p^N} = 0$ for every element $x$ of an ideal with divided powers of an $A$-algebra.
\end{corollary}

\begin{proof}
The integer $p^N$ divides $(p^N)!$, and the second statement is Lemma~\ref{lem:pd-exponent} for $n = p^N$.
\end{proof}

\leanref{\lean{natCast_factorial_pow_eq_zero}, \lean{dividedPowers_bounded}, in \lean{InfinitesimalSite/Envelope.lean}.}

The bound $n$ of Lemma~\ref{lem:pd-exponent} depends only on $A$, not on $D$ or $\delta$. So, if $n! = 0$ in $A$, an $A$-algebra map from $B$ to an $A$-algebra with divided powers that sends $\ker(B \twoheadrightarrow C)$ into the divided power ideal kills the $n$-th powers of the elements of $\ker(B \twoheadrightarrow C)$. Truncating $B$ by these powers therefore does not change its divided power envelope (Lemma~\ref{lem:truncation-envelope}), and the cohomology of $\Ocris$ is computed on the single site of thickenings of exponent $n$ (Theorem~\ref{thm:adjunction-derived}). We use $n = p^e$, for which the condition $(p^e)! = 0$ holds as soon as $p$ is nilpotent enough in $A$ (Remark~\ref{rem:valuation}). This uniformity is why we ask for a bounded exponent and not merely for a nil kernel.

\subsection{Comparison with the usual crystalline site}\label{subsec:comparison}

In Berthelot's crystalline site \cite{berthelot,berthelot-ogus}, objects are divided power thickenings $U \hookrightarrow T$ with $U$ open in $X = \Spec(C)$, the topology is the Zariski (or étale) topology, and the structure sheaf sends such an object to $\mathcal{O}_T$. For affine $X$ its cohomology is the derived limit of the tautological diagram $(B \twoheadrightarrow C) \mapsto B$ over the category of affine divided power thickenings of $C$ \cite[Tags~07HM, 07JK]{stacks}. In $\INF(C/A)$ every object is a thickening of all of $\Spec(C)$, objects carry no divided powers, and the topology is chaotic.

\spa

The second difference is the value of the structure sheaf. On an object $\pi \colon B \twoheadrightarrow C$ of the infinitesimal site, $\Ocris$ is a compatible divided power envelope of $(B, \ker(\pi))$. That $R\Gamma(\INF(C/A), \Ocris)$ agrees with the cohomology of Berthelot's site is classical: both are computed by the de Rham complex of a divided power envelope of a polynomial presentation (Theorem~\ref{thm:main} and \cite[Tag~07LG]{stacks}), and alternatively both are derived limits that can be compared through the envelope functor, which is left adjoint to the functor forgetting divided powers \cite[Tag~07JK]{stacks}. This agreement is not formalised (Section~\ref{sec:scope}).

%======================================================================
\section{The bounded-exponent infinitesimal site}\label{sec:site}
%======================================================================

\subsection{Objects and morphisms}

We write $\INF(C/A)$ for the category of Definition~\ref{def:thickening}, and call it the \emph{big infinitesimal site} of $C$ over $A$: every object is a thickening of $\Spec(C)$ itself, with arbitrary $B$, and the upper case follows Berthelot's notation for big sites. In Lean it is \lean{InfinitesimalThickening A C}, in \lean{InfinitesimalSite/Basic.lean}. The fields of an object are a carrier ring with an $A$-algebra structure, the map \texttt{$\pi$}, the field \lean{surjective}, and the field \lean{kernel_bounded}; no polynomial presentation is chosen or stored. A morphism is \lean{InfinitesimalThickening.Hom}, with fields \lean{toAlgHom} and \lean{overC}. A morphism $f$ automatically sends $\ker(\pi)$ into $\ker(\pi')$, so there is no further compatibility to check.

\subsection{Exponents and the Frobenius levels \texorpdfstring{$\INF_e$}{INF\_e}}

\begin{definition}
Let $X = (B, \pi)$ be an object of $\INF(C/A)$ and $m \in \N$. The object $X$ \emph{has exponent} $m$ if $x^m = 0$ for every $x \in \ker(\pi)$. The \emph{Frobenius exponent} $\exp(X)$ is the least $e \in \N$ such that $X$ has exponent $p^e$, that is, such that $x \mapsto x^{p^e}$ kills $\ker(\pi)$. For $e \in \N$, the \emph{Frobenius level} $\INF_e \subseteq \INF(C/A)$ is the full subcategory of objects having exponent $p^e$.
\end{definition}

Every object has some exponent $m$, hence the exponent $p^m \ge m$, so $\exp(X)$ is well defined. In Lean, \lean{Level A C m} is the full subcategory of objects of exponent $m$, for an arbitrary $m \in \N$, and $\INF_e$ is \lean{Level A C (p ^ e)}. The Frobenius exponent is defined for every integer $p > 1$.

\begin{proposition}\label{prop:levels}
Let $X$ be an object of $\INF(C/A)$ and let $e \le e'$ be integers.
\begin{itemize}
\item[\normalfont{(1)}] If $X$ has exponent $p^e$, then it has exponent $p^{e'}$; in particular $\INF_e \subseteq \INF_{e'}$.
\item[\normalfont{(2)}] The object $X$ lies in $\INF_e$ if and only if $\exp(X) \le e$.
\item[\normalfont{(3)}] The object $X$ lies in $\INF_{\exp(X)}$, so $\INF(C/A)$ is the increasing union of the Frobenius levels.
\end{itemize}
\end{proposition}

\leanref{\lean{HasExponent}, \lean{HasExponent.mono}, \lean{frobeniusExponent} (the Frobenius exponent $\exp(X)$, defined by \lean{Nat.find}), \lean{frobeniusExponent_le_iff}, \lean{hasExponent_pow_frobeniusExponent}, \lean{Level}, \texttt{Level.$\iota$}, \lean{Level.exponentLE_pow_mono}, \lean{Level.inclPow}, \lean{Level.ofObjFrobenius}, \lean{Level.exists_frobenius_level}. For arbitrary exponents: \lean{exponent}, the least $m$ such that $X$ has exponent $m$, with \lean{exponent_le_iff}, \lean{Level.incl}, \lean{Level.exists_level}, and \lean{exponent_le_pow_frobeniusExponent}, which bounds it by $p^{\exp(X)}$.}

\subsection{Truncation}

For an ideal $K$ of a ring and $m \in \N$, let $\pow(K, m)$ be the ideal generated by $\{x^m : x \in K\}$.

\begin{lemma}\label{lem:powideal}
Let $K$ be an ideal of a ring and let $m > 0$. Then $\pow(K,m) \subseteq K$. If $f$ is a ring map with $f(x)^m = 0$ for all $x \in K$, then $\pow(K,m) \subseteq \ker(f)$.
\end{lemma}

\leanref{\lean{powIdeal}, \lean{pow_mem_powIdeal}, \lean{powIdeal_le}, \lean{powIdeal_le_ker}.}

Let $\pi \colon B \twoheadrightarrow C$ be a surjection of $A$-algebras and let $e \in \N$. We truncate $B$ by the ideal generated by the $p^e$-th powers of elements of the kernel:
\[
T_e(B) \;\coloneqq\; B \,\big/\, \pow(\ker(\pi), p^e) \;=\; B \,\big/\, \big( x^{p^e} \;:\; x \in \ker(\pi) \big).
\]

\begin{proposition}\label{prop:truncation}
Let $\pi \colon B \twoheadrightarrow C$ be a surjection of $A$-algebras and let $e \in \N$. The kernel of the induced map $T_e(B) \twoheadrightarrow C$ is the image of $\ker(\pi)$, and $T_e(B) \twoheadrightarrow C$ is an object of $\INF_e$. Every $A$-algebra map $f \colon B \to Y$ over $C$, with $Y$ in $\INF_e$, factors uniquely through $B \twoheadrightarrow T_e(B)$.
\end{proposition}

\begin{proof}
Since $p^e > 0$, $\pow(\ker(\pi), p^e) \subseteq \ker(\pi)$ by Lemma~\ref{lem:powideal}, so $\pi$ factors through $T_e(B)$, and as $B \twoheadrightarrow T_e(B)$ is surjective the kernel of $T_e(B) \twoheadrightarrow C$ is the image of $\ker(\pi)$; the $p^e$-th power of the class of $x \in \ker(\pi)$ is the class of $x^{p^e} \in \pow(\ker(\pi), p^e)$, hence zero.

Since $\pi_Y \circ f = \pi$, the map $f$ sends $\ker(\pi)$ into $\ker(\pi_Y)$, so it kills $p^e$-th powers of elements of $\ker(\pi)$ and, by Lemma~\ref{lem:powideal}, the ideal $\pow(\ker(\pi), p^e)$. It therefore factors through the quotient, and the factorisation is unique because $B \twoheadrightarrow T_e(B)$ is surjective.
\end{proof}

\leanref{\lean{InfinitesimalThickening.truncate}, \lean{truncate_hasExponent}, \lean{truncateLift}, \lean{truncateLift_comp_mk}, \lean{truncate_hom_ext}, \lean{ker_truncate_eq_map}; the truncation is defined by $m$-th powers for every $m > 0$, and $T_e$ is the case $m = p^e$.}

The kernels we need are in general not finitely generated: for the kernel of $A[x_c : c \in C] \twoheadrightarrow C$, there is one variable for each element of $C$. For such kernels a bounded exponent is strictly weaker than nilpotence of the ideal.

\begin{example}
Consider the surjection $\F_p[y_1, y_2, \dots]/(y_1^p, y_2^p, \dots) \twoheadrightarrow \F_p$ sending every $y_i$ to $0$. Every element of the kernel has vanishing $p$-th power, because the Frobenius is additive and every monomial in the kernel is divisible by some $y_i$, so this object lies in $\INF_1$. The kernel is not nilpotent, since $y_1 \cdots y_n \ne 0$ lies in its $n$-th power for every $n$.
\end{example}

Divided power ideals behave in the same way: when $(p^e)! = 0$ in $A$, every element of a divided power ideal of an $A$-algebra has vanishing $p^e$-th power by Lemma~\ref{lem:pd-exponent}, while the ideal need not be nilpotent. This is why the truncation uses the ideal of $p^e$-th powers and not a power of the kernel.

For an object $X = (B, \pi)$ of $\INF(C/A)$ we write $T_e(X)$ for the object $T_e(B) \twoheadrightarrow C$ of $\INF_e$, and $u_X \colon X \twoheadrightarrow T_e(X)$ for the quotient map $B \twoheadrightarrow T_e(B)$, a morphism of $\INF(C/A)$. Let $\iota_e \colon \INF_e \hookrightarrow \INF(C/A)$ be the inclusion.

\begin{theorem}\label{thm:adjunction}
Let $e \in \N$. The assignment $X \mapsto T_e(X)$ extends to a functor $T_e \colon \INF(C/A) \to \INF_e$, which is left adjoint to $\iota_e$. The unit of the adjunction at $X$ is $u_X$.
\end{theorem}

\begin{proof}
For a morphism $f \colon X \to Y$, Proposition~\ref{prop:truncation} applied to $u_Y \circ f$ gives a unique $T_e(f) \colon T_e(X) \to T_e(Y)$ with $T_e(f) \circ u_X = u_Y \circ f$; functoriality follows from uniqueness. For $Y$ in $\INF_e$, Proposition~\ref{prop:truncation} says that composition with $u_X$ is a bijection
\[
\Hom_{\INF_e}\big(T_e(X), Y\big) \;\simeq\; \Hom_{\INF(C/A)}\big(X, \iota_e(Y)\big),
\]
and this bijection is natural in $X$ and $Y$ because $T_e(f) \circ u_X = u_Y \circ f$.
\end{proof}

\leanref{\lean{Level.truncation}, \lean{InfinitesimalThickening.toTruncate}, \lean{Level.toTruncate_comp_truncation_map}, \lean{Level.truncation_hom_ext}, \lean{Level.truncationHomEquiv}, \lean{Level.truncationAdjunction}, \lean{Level.truncationAdjunction_unit_app}, in \lean{InfinitesimalSite/Filtration.lean}. The inclusion $\iota_e$ is \texttt{Level.$\iota$}.}

\subsection{Binary coproducts}

\begin{proposition}\label{prop:coproducts}
Let $e \in \N$ and let $X = (B, \pi)$ and $X' = (B', \pi')$ be objects of $\INF_e$. The truncation of the map $B \otimes_A B' \twoheadrightarrow C$ sending $b \otimes b'$ to $\pi(b)\,\pi'(b')$ is a coproduct of $X$ and $X'$ in $\INF_e$. In particular $\INF_e$ has binary coproducts, i.e.\ binary products in the site orientation $\INF_e^{\mathrm{op}}$.
\end{proposition}

\begin{proof}
For maps $f \colon X \to Y$ and $g \colon X' \to Y$, the map out of the coproduct is the factorisation given by Proposition~\ref{prop:truncation} of $b \otimes b' \mapsto f(b) g(b')$. Uniqueness follows from $b \otimes b' = (b \otimes 1)(1 \otimes b')$.
\end{proof}

\leanref{\lean{Level.coprod}, \lean{Level.inl}, \lean{Level.inr}, \lean{Level.coprodDesc}, \lean{Level.coprod_hom_ext}, \lean{Level.coprodCofanIsColimit}, instance \lean{Level.hasBinaryCoproducts} for \lean{Level A C m} under \lean{NeZero m}, applied to $m = p^e$.}

\subsection{The weakly initial object and the Čech nerve}

For an object $Y = (B', \pi')$ let $S(Y)$ be the set of maps of sets $s \colon C \hookrightarrow B'$ with $\pi' \circ s = \mathrm{id}$. It is nonempty because $\pi'$ is surjective, and it is functorial in $Y$. Write $[n] \coloneqq \{0, \dots, n\}$. For $n \ge 0$ let $X_e(n)$ be the truncation $T_e$ of the map $A[x_{i,c} : i \in [n],\ c \in C] \twoheadrightarrow C$ sending $x_{i,c}$ to $c$, that is, the quotient of the polynomial ring by the $p^e$-th powers of the elements of the kernel.

\leanref{\lean{InfinitesimalThickening.Sections}, \lean{sections_nonempty}, \lean{Sections.map}, \lean{Level.cechObj}; the object $X_e(n)$ is \lean{Level.cechObj (p ^ e) n}.}

A map $\theta \colon [m] \to [n]$ induces the morphism $X_e(m) \to X_e(n)$ sending $x_{i,c}$ to $x_{\theta(i),c}$ (\lean{cechMap}).

\begin{proposition}\label{prop:cech-hom}
Let $e, n \ge 0$ and let $Y$ be an object of $\INF_e$. Maps $X_e(n) \to Y$ correspond bijectively to $(n+1)$-tuples $(s_0, \dots, s_n)$ of elements of $S(Y)$, via $x_{i,c} \mapsto s_i(c)$. The bijection is natural in $Y$, and precomposition with the map induced by $\theta \colon [m] \to [n]$ corresponds to $s \mapsto s \circ \theta$. In particular $X_e(0)$ maps to every object of $\INF_e$.
\end{proposition}

\leanref{\lean{Level.cechHomEquiv}, \lean{cechHomEquiv_cechMap}, \lean{cechHomEquiv_comp}, \lean{Level.nonempty_hom_cechObj_zero}.}

Restricted to monotone maps, these renamings of variables give a cosimplicial object $X_e(\bullet)$ of $\INF_e$ (\lean{Level.cechNerve}, \lean{cechMap}, \lean{cechMap_id}, \lean{cechMap_comp}). In the site orientation, $X_e(0)$ is weakly final, $\INF_e$ has binary products, and $X_e(\bullet)$ is the Čech nerve of $X_e(0)$ by Proposition~\ref{prop:cech-hom}; these are the hypotheses of \cite[Tag~07JM]{stacks}. The polynomial rings are only a device in proofs: they enter neither the definition of the site nor that of $\Ocris$ or of $R\Gamma(\INF(C/A), \Ocris)$; they appear only through the objects $X_e(n)$ in the statements of Theorems~A and~B.

\begin{remark}\label{rem:general-exponent}
Proposition~\ref{prop:truncation}, Theorem~\ref{thm:adjunction} and Propositions~\ref{prop:coproducts} and~\ref{prop:cech-hom} only use that $p^e$ is a positive integer. For every integer $m > 0$ they hold, with the same proofs, for the full subcategory of thickenings of exponent $m$, with the truncation by $m$-th powers in place of $T_e$. This is the generality of the Lean statements, which concern \lean{Level A C m} under \lean{NeZero m}; the Frobenius level $\INF_e$ is the case $m = p^e$.
\end{remark}

\subsection{The topology}

The site $\INF(C/A)$ carries the chaotic topology, i.e.\ Mathlib's bottom Grothendieck topology on the opposite category; we use the chaotic topology on $\INF_e$ as well. Every presheaf is a sheaf for it; for $\Ocris$ this is \lean{isOcrisSheaf_chaotic}, with the topology packaged as \lean{TopologyData.chaotic}. For the chaotic topology, sheaf cohomology is the right derived functor of the limit, and we take this as the definition of $R\Gamma$ (Section~\ref{subsec:derived}); the Lean statements use the derived limit directly. This does not make higher cohomology vanish.

%======================================================================
\section{The structure sheaf \texorpdfstring{$\Ocris$}{Ocris}}\label{sec:ocris}
%======================================================================

From now on we work in the following setting.

\begin{setting}\label{Setting:s}
Let $p$ be a prime, $A$ a ring and $I \subseteq A$ an ideal with divided powers $\gamma$, such that $p \in I$ and $p$ is nilpotent in $A$. Let $e$ be a positive integer with $(p^e)! = 0$ in $A$, and let $C$ be an $A$-algebra with $I \cdot C = 0$.
\end{setting}

\begin{remark}\label{rem:valuation}
Since $p$ is nilpotent in $A$, every integer prime to $p$ is a unit in $A$, so the prime-to-$p$ part of $(p^e)!$ is a unit. By Legendre's formula, the $p$-adic valuation of $(p^e)!$ is $1 + p + \dots + p^{e-1} = (p^e-1)/(p-1)$. Hence the hypothesis $(p^e)! = 0$ holds if and only if $p^{(p^e-1)/(p-1)} = 0$ in $A$. If $p^N = 0$ in $A$, it holds for $e = N$ by Corollary~\ref{cor:pd-exponent-nilpotent}, and already for the least $e$ with $(p^e-1)/(p-1) \ge N$, that is, with $p^e \ge (p-1)N + 1$. For this $e$ the exponent $p^e$ of the objects of $\INF_e$ is smaller than $p\big((p-1)N + 1\big)$, which is linear in $N$, whereas the choice $e = N$ gives the exponent $p^N$. This remark is not formalised; the library uses only Corollary~\ref{cor:pd-exponent-nilpotent}.
\end{remark}

In Lean the conditions on $p$ form the class \lean{IsCrystallineBase p A I}, the condition on $e$ is the hypothesis \texttt{he : ((Nat.factorial (p \textasciicircum{} e) : $\mathbb{N}$) : A) = 0}, and $I \cdot C = 0$ is the hypothesis \lean{hIC : I.map (algebraMap A C)}~$= \bot$. We refer to Setting~\ref{Setting:s} as the Basic Setting. For an object $X = (B, \pi)$ of $\INF(C/A)$ we write $J_X \coloneqq \ker(\pi)$.

\subsection{Envelopes as chosen data}

\begin{definition}\label{def:envelope}
Let $B$ be an $A$-algebra and $J \subseteq B$ an ideal. A \emph{compatible divided power envelope} of $(B,J)$ is an $A$-algebra $D$ with a map $\varphi \colon B \to D$, an ideal $K \supseteq JD + ID$ and divided powers $\delta$ on $K$ compatible with $\gamma$, that is, such that $A \to D$ is a morphism of divided power rings, with the following property. For every $A$-algebra $D'$, every $A$-algebra map $\varphi' \colon B \to D'$ and all divided powers $\delta'$ on an ideal $K' \supseteq \varphi'(J)D' + ID'$ compatible with $\gamma$, there is a unique morphism of divided power algebras $\psi \colon D \to D'$ with $\psi \circ \varphi = \varphi'$.
\end{definition}

This is the predicate \lean{IsDPEnvelope}, defined in \lean{PDDifferentials.lean}; \lean{DPEnvelopeChoice} bundles the data and \lean{HasDPEnvelope} is the propositional existence class. The ideal $K$ is part of the data. One can show that it is the divided power ideal generated by $J$ and $I$; this is neither formalised nor used. It is in general strictly larger than $JD + ID$. For example, let $A = \mathbb{Z}/p^N$ with $I = pA$ and its canonical divided powers, $B = A[x]$ and $J = (x)$. The envelope is $A\langle x \rangle$, and $x^{[p]} \in K$ but $x^{[p]} \notin (x,p)A\langle x \rangle$: since $x \cdot x^{[k]} = (k+1)\,x^{[k+1]}$, every element of $(x,p)A\langle x\rangle$ has its coefficient of $x^{[p]}$ in $pA$.

\spa

For an object $X$ of $\INF(C/A)$, \lean{EnvelopeChoice X} (in \lean{InfinitesimalSite/Envelope.lean}) packages an envelope $(D_X, \varphi_X, K_X, \delta_X)$ of $(B, J_X)$. An \emph{envelope family} $E$ (\lean{EnvelopeFamily}) is a choice of such an envelope for every object of $\INF(C/A)$. Every result about $\Ocris$ is stated for an arbitrary envelope family. By Theorem~\ref{thm:envelope-existence} an envelope family exists (\lean{envelopeExistenceData}, \lean{envelopeFamily} in \lean{MainTheorems.lean}).

\subsection{Existence of envelopes}

\begin{theorem}\label{thm:envelope-existence}
Let $(A, I, \gamma)$ be a divided power ring, $B$ an $A$-algebra and $J \subseteq B$ an ideal. Then $(B, J)$ has a compatible divided power envelope.
\end{theorem}

\begin{proof}
We follow the representability argument of \cite[Tags~07H8, 07GV, 07GW]{stacks}, with an explicit indexing set; there $IB \subseteq J$ is assumed, and Definition~\ref{def:envelope} is the envelope of $J + IB$. Let $R \coloneqq B[X_{n,f} : n \ge 0,\ f \in J]$; the variable $X_{n,f}$ stands for $\delta_{n+1}(f)$. For an ideal $\mathfrak{a} \subseteq R$, let $K(\mathfrak{a}) \subseteq R/\mathfrak{a}$ be the ideal generated by the classes of the variables and by the image of $I$. Call a pair $(\mathfrak{a}, \delta)$, with $\delta$ divided powers on $K(\mathfrak{a})$, \emph{admissible} if the class of every $f \in J$ lies in $K(\mathfrak{a})$, $X_{n,f} = \delta_{n+1}(f)$ in $R/\mathfrak{a}$, and $A \to R/\mathfrak{a}$ is a morphism of divided power rings. Admissible pairs form a set $S$. Let $\mathfrak{b}$ be the kernel of $R \to \prod_{s \in S} R/\mathfrak{a}_s$ and $D \coloneqq R/\mathfrak{b}$. The product of the ideals $K(\mathfrak{a}_s)$ carries the componentwise divided powers, and the ideal $K(\mathfrak{b})$ of $D \hookrightarrow \prod_s R/\mathfrak{a}_s$ is stable under them. Indeed, it suffices to check this on generators, and for $m \ge 1$ we have $\delta_m(X_{n,f}) = \frac{(m(n+1))!}{m!\,((n+1)!)^m}\, X_{m(n+1)-1,f}$ (because $X_{n,f} = \delta_{n+1}(f)$ in each factor) and $\delta_m(a) = \gamma_m(a)$ for $a \in I$, uniformly in $s$. So $K(\mathfrak{b})$ has divided powers, compatible with $\gamma$ and containing $JD + ID$ (the class of $f$ equals that of $X_{0,f}$ in every factor).

Let $(D', K', \delta')$ and $\varphi' \colon B \to D'$ be as in Definition~\ref{def:envelope}. Let $\mathfrak{a}$ be the kernel of the evaluation $R \to D'$ sending $X_{n,f}$ to $\delta'_{n+1}(\varphi'(f))$. Then $R/\mathfrak{a} \hookrightarrow D'$ and $K(\mathfrak{a}) \subseteq K'$, since $\varphi'(J) \subseteq K'$ and $ID' \subseteq K'$. The ideal $K(\mathfrak{a})$ is stable under $\delta'$ by the same computation, the class of $f \in J$ equals that of $X_{0,f}$ because $\delta'_1(\varphi'(f)) = \varphi'(f)$, and so the restriction of $\delta'$ makes $\mathfrak{a}$ admissible. The projection $D \twoheadrightarrow R/\mathfrak{a} \hookrightarrow D'$ is the required morphism. It is unique because $D$ is generated over $B$ by the classes of the $X_{n,f} = \delta_{n+1}(f)$.
\end{proof}

\leanref{\lean{DPEnvelope.hasDPEnvelope}, \lean{DPEnvelope.env_isDPEnvelope}, \lean{DPEnvelope.existsUnique_lift}, \lean{DPEnvelope.IsAdmissible}, \lean{DPEnvelope.envDividedPowers}, \lean{DPEnvelope.envPDIdeal_stab}, \lean{piDividedPowers}, \lean{restrictDividedPowers}, \lean{exists_dpow_mem_span}, in \lean{Envelope/Construction.lean}.}

The proof gives no explicit description of $D$, but none is needed: every other result uses envelopes only through their universal property.

\subsection{Consequences of the universal property}

\begin{proposition}\label{prop:envelope-unique}
Two compatible divided power envelopes of the same $(B, J)$ are isomorphic as $A$-algebras, by mutually inverse comparison maps given by the universal property.
\end{proposition}

\leanref{\lean{IsDPEnvelope.comparison}, \lean{comparison_spec}, \lean{identity_spec}, \lean{comp_spec}, \lean{IsDPEnvelope.comparisonEquiv}, \lean{DPEnvelopeChoice.hom_ext}.}

\begin{proposition}\label{prop:functoriality}
Let $E$ be an envelope family. For a morphism $f \colon X \to Y$ of $\INF(C/A)$ there is a unique morphism of divided power algebras $c(f) \colon D_X \to D_Y$ with $c(f) \circ \varphi_X = \varphi_Y \circ f$. Moreover $c(\mathrm{id}_X) = \mathrm{id}$ and $c(g \circ f) = c(g) \circ c(f)$, and for two choices of envelope at the same $X$ the maps $c(\mathrm{id}_X)$ in the two directions are mutually inverse.
\end{proposition}

\begin{proof}
Apply the universal property of $D_X$ to $\varphi_Y \circ f$; this is possible because $f(J_X) \subseteq J_Y$. The identities follow from uniqueness.
\end{proof}

\leanref{\lean{EnvelopeChoice.comparison}, \lean{comparison_spec}, \lean{comparison_id}, \lean{comparison_comp}, \lean{comparison_comp_comparison_left}, \lean{comparison_comp_comparison_right}.}

\begin{corollary}\label{cor:ocris}
The assignment $X \mapsto D_X$ is a functor from $\INF(C/A)$ to commutative $A$-algebras, i.e.\ a presheaf of rings $\Ocris$ on the site. Two envelope families give isomorphic presheaves.
\end{corollary}

\leanref{\lean{ocris}, \lean{ocrisCommRing}, \lean{ocris_map_spec}, \lean{ocrisOfExistenceData}, \lean{envelopeFamilyIso}, \lean{ocrisFamilyIso}.}

No proof looks inside $D_X$; everything goes through \lean{IsDPEnvelope}.

\subsection{The projection to \texorpdfstring{$C$}{C}}

The next result identifies the divided power ideal of an envelope with a kernel.

\begin{theorem}\label{thm:projection}
Assume the Basic Setting, let $E$ be an envelope family and let $X = (B, \pi)$ be an object of $\INF(C/A)$. There is a unique morphism of divided power algebras $q \colon D_X \to C$, for the zero divided power ideal of $C$, such that $q \circ \varphi_X = \pi$. The map $q$ is surjective and
\[
\ker(q) \;=\; K_X .
\]
In particular $x^{p^e} = 0$ for every $x \in \ker(q)$. If $q'$ is the corresponding map for an object $Y$ and $f \colon X \to Y$ is a morphism, then $q' \circ c(f) = q$.
\end{theorem}

\begin{proof}
Existence and uniqueness of $q$ follow from the universal property, since $I\cdot C = 0$ and $\pi(J_X) = 0$. Since $q$ is a morphism of divided power algebras into $(C, 0)$, we have $K_X \subseteq \ker(q)$. Conversely, since $J_X D_X + I D_X \subseteq K_X$ and $\pi$ is surjective, $D_X \twoheadrightarrow D_X/K_X$ induces a map $i \colon C \to D_X/K_X$ with $i \circ \pi = (D_X \twoheadrightarrow D_X/K_X) \circ \varphi_X$. The projection $D_X \twoheadrightarrow D_X/K_X$ and $i \circ q$ are then two morphisms of divided power algebras into $(D_X/K_X, 0)$ solving the same universal problem, so they agree; hence $\ker(q) \subseteq K_X$. The exponent bound is Lemma~\ref{lem:pd-exponent} for $n = p^e$. Finally, $q' \circ c(f)$ and $q$ are morphisms of divided power algebras with $q' \circ c(f) \circ \varphi_X = q' \circ \varphi_Y \circ f = \pi$, so they agree by uniqueness.
\end{proof}

\leanref{\lean{EnvelopeChoice.projection}, \lean{projection_spec}, \lean{existsUnique_projection}, \lean{projection_surjective}, \lean{projectionKernel_eq_dividedPowersIdeal}, \lean{dividedPowers_pow_eq_zero_of_factorial_eq_zero}, \lean{dividedPowers_bounded}, \lean{comparison_overC}.}

\subsection{Envelopes of truncations}

The following lemma is the only place where the bound of Lemma~\ref{lem:pd-exponent} enters the comparison of $\INF(C/A)$ with $\INF_e$. Recall the quotient maps $u_X \colon X \twoheadrightarrow T_e(X)$ of Theorem~\ref{thm:adjunction}.

\begin{lemma}\label{lem:truncation-envelope}
Assume the Basic Setting.
\begin{itemize}
\item[\normalfont{(1)}] Let $\pi \colon B \twoheadrightarrow C$ be a surjection of $A$-algebras. A compatible divided power envelope of $T_e(B)$, with respect to the kernel of $T_e(B) \twoheadrightarrow C$ and composed with $B \twoheadrightarrow T_e(B)$, satisfies the universal property of a compatible divided power envelope of $(B, \ker(\pi))$.
\item[\normalfont{(2)}] Let $E$ be an envelope family. For every object $X$ of $\INF(C/A)$ the map $c(u_X) \colon D_X \to D_{T_e(X)}$ of Proposition~\ref{prop:functoriality} is bijective, and these maps form an isomorphism of functors $\Ocris \simeq \Ocris \circ \iota_e \circ T_e$ on $\INF(C/A)$.
\end{itemize}
\end{lemma}

\begin{proof}
For (1), let $\varphi' \colon B \to D'$ be a map sending $J \coloneqq \ker(\pi)$ into an ideal $K'$ with divided powers compatible with $\gamma$. For $x \in J$, Lemma~\ref{lem:pd-exponent} for $n = p^e$ gives $\varphi'(x)^{p^e} = 0$, so $\varphi'$ kills $\pow(J, p^e)$ by Lemma~\ref{lem:powideal} and factors through $T_e(B)$, whose kernel is the image of $J$ by Proposition~\ref{prop:truncation}. The universal property of the envelope of $T_e(B)$ concludes.

\spa

For (2), write $X = (B, \pi)$. By (1), applied to $\varphi_X$, there is a morphism of divided power algebras $\psi \colon D_{T_e(X)} \to D_X$ with $\psi \circ \varphi_{T_e(X)} \circ u_X = \varphi_X$. Then $\psi \circ c(u_X)$ is a morphism of divided power algebras with $\psi \circ c(u_X) \circ \varphi_X = \varphi_X$, hence the identity by the universal property of $D_X$. Similarly $c(u_X) \circ \psi \circ \varphi_{T_e(X)}$ and $\varphi_{T_e(X)}$ agree after composition with $u_X$, hence agree since $u_X$ is surjective, and $c(u_X) \circ \psi$ is the identity by the universal property of $D_{T_e(X)}$. Naturality in $X$ follows from Proposition~\ref{prop:functoriality} and the identity $T_e(f) \circ u_X = u_Y \circ f$ of Theorem~\ref{thm:adjunction}.
\end{proof}

\leanref{\lean{truncIdeal_le_ker_of_dividedPowers}, \lean{EnvelopeChoice.truncate_isInitial}, \lean{EnvelopeChoice.comparison_toTruncate_bijective}, \lean{EnvelopeFamily.isIso_ocrisModuleLarge_map_toTruncate}, \lean{EnvelopeFamily.ocrisTruncationIso}, in \lean{Ocris/Filtration.lean}. They are stated for the truncation by $m$-th powers, for every $m > 0$ with $m! = 0$ in $A$, and the lemma is the case $m = p^e$. Statement (1) is the existence-and-uniqueness clause of Definition~\ref{def:envelope}; the other two conditions follow from those of the envelope of $T_e(B)$.}

\subsection{Derived global sections}\label{subsec:derived}

The derived global sections $R\Gamma(\INF(C/A), \Ocris)$ are, by definition, the right derived functor of the limit applied to the diagram of $A$-modules underlying $\Ocris$. They are computed with injective resolutions in the functor category, which has enough injectives by Theorem~\ref{thm:enough-injectives}. In Lean this object of the homotopy category is \lean{ocrisDerivedGlobalSections E}; for universe reasons the diagram takes values in \lean{ModuleCat.{u+1} A}, through \lean{ULift}. Neither the derived nor the underived global sections depend on the envelope family (\lean{ocrisDerivedGlobalSectionsIso}, \lean{ocrisGlobalSectionsIso}).

\spa

For $e \in \N$, the derived global sections $R\Gamma(\INF_e, \Ocris)$ are defined in the same way from the restriction $\Ocris \circ \iota_e$ of this diagram to $\INF_e$ (\lean{levelOcrisModule}).

\begin{theorem}\label{thm:adjunction-derived}
Assume the Basic Setting and let $E$ be an envelope family. There is an isomorphism in the homotopy category of complexes of $A$-modules
\[
R\Gamma(\INF(C/A), \Ocris) \;\simeq\; R\Gamma(\INF_e, \Ocris).
\]
\end{theorem}

\begin{proof}
By Theorem~\ref{thm:adjunction}, the truncation $T_e \colon \INF(C/A) \to \INF_e$ is left adjoint to the inclusion $\iota_e$. Theorem~\ref{thm:derived-left-adjoint}, applied to $T_e$ and to $\Ocris \circ \iota_e$, gives an isomorphism
\[
R\lim{}_{\INF_e}(\Ocris \circ \iota_e) \;\simeq\; R\lim{}_{\INF(C/A)}(\Ocris \circ \iota_e \circ T_e).
\]
The left-hand side is $R\Gamma(\INF_e, \Ocris)$. By Lemma~\ref{lem:truncation-envelope}(2), $\Ocris \circ \iota_e \circ T_e$ is isomorphic to $\Ocris$, so the right-hand side is isomorphic to $R\lim_{\INF(C/A)}(\Ocris) = R\Gamma(\INF(C/A), \Ocris)$.
\end{proof}

\leanref{\lean{EnvelopeFamily.ocrisDerivedGlobalSectionsIsoLevel}, \lean{levelOcrisModule}, \lean{Level.truncationAdjunction}, \lean{Cohomology.derivedLimitIsoOfAdjunction}, \lean{EnvelopeFamily.ocrisTruncationIso}, in \lean{Ocris/Filtration.lean}. In Lean the isomorphism is proved on \lean{Level A C m} for every $m > 0$ with $m! = 0$ in $A$.}

The proof uses the single site $\INF_e$; no limit over $e$ is involved.

\subsection{The structure sheaf on the Čech objects}

For $n \ge 0$ let $P_n \coloneqq A[x_{i,c} : i \in [n],\ c \in C]$, with the map $P_n \twoheadrightarrow C$ sending $x_{i,c}$ to $c$, and let $D(n) \coloneqq \Ocris(X_e(n))$ be the envelope chosen by $E$ at $X_e(n)$.

\begin{proposition}\label{prop:cech-envelope}
Assume the Basic Setting. For every $n \ge 0$, the algebra $D(n)$, with the composite $P_n \twoheadrightarrow X_e(n) \to D(n)$, satisfies the universal property of the compatible divided power envelope of $P_n \twoheadrightarrow C$.
\end{proposition}

\begin{proof}
The object $X_e(n)$ is the truncation of $P_n \twoheadrightarrow C$, so this is Lemma~\ref{lem:truncation-envelope}(1).
\end{proof}

\leanref{\lean{ocrisCech_isEnvelope}, in \lean{Ocris/Filtration.lean}.}

Thus the Čech complex $n \mapsto \Ocris(X_e(n))$ is the complex $n \mapsto D(n)$ of envelopes of the polynomial presentations, as in \cite[Tag~07JM]{stacks}.

\subsection{The de Rham complex on the site}

Using the divided power de Rham complex of Section~\ref{sec:pd-derham}, the assignment $X \mapsto \Omega^\bullet_{D_X}$ is a functor from $\INF(C/A)$ to cochain complexes of $A$-modules, with a natural map of complexes $\Omega^\bullet_{D_X} \to \Ocris(X)[0]$ which is the identity in degree $0$. It is independent of the envelope family up to natural isomorphism.

\leanref{\lean{ocrisDeRhamComplex}, \lean{ocrisDeRhamAugmentation}, \lean{ocrisDeRhamComplexIso}, \lean{ocrisDeRhamZeroIso}.}

%======================================================================
\section{Divided power differentials and de Rham complexes}\label{sec:pd-derham}
%======================================================================

In this section $B$ is an $A$-algebra and $J \subseteq B$ an ideal with divided powers $\delta$; we write $x\dpow{n} \coloneqq \delta_n(x)$ for $x \in J$. The results follow \cite[Tag~07HQ]{stacks}.

\subsection{Divided power Kähler differentials}

\begin{definition}
A \emph{PD derivation} of $B$ into a $B$-module $M$ is an $A$-linear map $\theta \colon B \to M$ with $\theta(bb') = b\theta(b') + b'\theta(b)$ for $b, b' \in B$ and $\theta(x\dpow{n}) = x\dpow{n-1}\theta(x)$ for $n \ge 1$ and $x \in J$. The module of \emph{PD Kähler differentials} is
\[
\Omega_{B,\delta} \;\coloneqq\; \Omega_{B/A} \big/ \big\langle d(x\dpow{n}) - x\dpow{n-1}\,dx \;:\; n \ge 1,\ x \in J \big\rangle_B ,
\]
with $d \colon B \to \Omega_{B,\delta}$ the induced map. For $q \ge 0$ put $\Omega^q_{B,\delta} \coloneqq \bigwedge^q_B \Omega_{B,\delta}$.
\end{definition}

\begin{proposition}\label{prop:pd-kaehler}
Let $B$, $J$ and $\delta$ be as above.
\begin{itemize}
\item[\normalfont{(1)}] The map $d \colon B \to \Omega_{B,\delta}$ is a PD derivation, and for every $B$-module $M$ composition with $d$ is a bijection from $\Hom_B(\Omega_{B,\delta}, M)$ to the set of PD derivations $B \to M$.
\item[\normalfont{(2)}] The $B$-module $\Omega_{B,\delta}$ is generated by $dB$.
\item[\normalfont{(3)}] A morphism of divided power algebras $f \colon (B,\delta) \to (D,\delta')$ induces an $f$-semilinear map $\Omega_{B,\delta} \to \Omega_{D,\delta'}$ with $db \mapsto d f(b)$.
\end{itemize}
\end{proposition}

\leanref{\lean{IsPDDerivation}, \lean{pdRelations}, \lean{PDKaehler}, \lean{pdD}, \lean{pdD_isPDDerivation}, \lean{PDKaehler.lift}, \lean{PDKaehler.lift_comp}, \lean{PDKaehler.lift_unique}, \lean{pdKaehlerDerivationEquiv}, \lean{span_range_pdD}, \lean{DeRham.pdKaehlerMap}, \lean{pdKaehlerMap_pdD}, \lean{pdKaehlerMap_smul}.}

\subsection{Square-zero extensions and derivations of envelopes}

\begin{proposition}\label{prop:square-zero}
Let $M$ be an $A$-module. The ideal $I \oplus M$ of the square-zero extension $A \oplus M$ carries the divided powers given, for $n \ge 1$, by
\[
\delta_n(x, z) \;\coloneqq\; \big(\gamma_n(x),\ \gamma_{n-1}(x)\,z\big),
\]
and $A \hookrightarrow A \oplus M$ is a morphism of divided power algebras \cite[Tag~07HH]{stacks}. If $\theta \colon D \to M$ is a PD derivation of a divided power algebra $D$, then $b \mapsto (b, \theta(b))$ and the projection $D \oplus M \twoheadrightarrow D$ are morphisms of divided power algebras.
\end{proposition}

\leanref{\lean{sqZeroDividedPowers}, \lean{sqZero_isDPMorphism} (in \lean{SquareZero.lean}; the composition axiom uses a \lean{Nat.uniformBell} identity), \lean{derivationGraph}, \lean{pdDerivationGraph_isDPMorphism}, \lean{sqZero_fst_isDPMorphism}.}

\begin{theorem}\label{thm:derivation-extension}
Let $\varphi \colon B \to D$ be an $A$-algebra map and $K$ an ideal of $D$ with divided powers $\delta$ compatible with $\gamma$, $JD + ID \subseteq K$, such that $\varphi$ is initial among maps to compatible divided power algebras sending $J$ into the divided power ideal; for instance, a compatible divided power envelope of $(B, J)$. Here $J$ carries no divided powers. Then every $A$-derivation $v \colon B \to M$ into a $D$-module extends uniquely to a PD derivation $\theta \colon D \to M$ with $\theta \circ \varphi = v$.
\end{theorem}

\begin{proof}
Give $D \oplus M$ the divided powers of Proposition~\ref{prop:square-zero} and apply initiality to the map $B \to D \oplus M$ sending $b$ to $(\varphi(b), v(b))$.
The first component of the resulting morphism $\psi \colon D \to D \oplus M$ is $\mathrm{id}_D$ by uniqueness, and its second component is $\theta$; $\psi(x\dpow{n}) = \psi(x)\dpow{n} = (x\dpow{n}, x\dpow{n-1}\theta(x))$ shows that $\theta$ is a PD derivation. For uniqueness, the graph of any other extension is another such morphism.
\end{proof}

\leanref{\lean{existsUnique_pdDerivation_extension} (for \lean{IsDPEnvelope}), \lean{existsUnique_pdDerivation_of_initial} (general form, in \lean{Ocris/PDDerivationExtension.lean}).}

The following is \cite[Tag~07HW]{stacks}.

\begin{theorem}\label{thm:envelope-differentials}
Let $D$ be a compatible divided power envelope of $(B, J)$, with divided powers $\delta$ on $K \subseteq D$. The canonical map $D \otimes_B \Omega_{B/A} \to \Omega_{D,\delta}$ sending $d' \otimes db$ to $d'\, d\varphi(b)$ is an isomorphism of $D$-modules.
\end{theorem}

\begin{proof}
Let $\theta \colon D \to D \otimes_B \Omega_{B/A}$ be the PD derivation extending $b \mapsto 1 \otimes db$ (Theorem~\ref{thm:derivation-extension}). Its factorisation through $\Omega_{D,\delta}$ (Proposition~\ref{prop:pd-kaehler}) is the inverse. One composite is the identity on the generators $d' \otimes db$; for the other, both composites with $d$ are PD derivations extending $b \mapsto db$, so they agree by uniqueness.
\end{proof}

\leanref{\lean{pdKaehlerDpEnvelopeComparison}, \lean{pdKaehler_dpEnvelope}, \lean{pdKaehlerDpEnvelopeEquiv}, in \lean{PDDifferentials.lean} and \lean{Envelope/PDKaehlerComparison.lean}.}

\subsection{Exterior powers of a quotient}

The following lemma of linear algebra is needed to descend the de Rham differential. It is proved in the library because we did not find it in Mathlib.

\begin{lemma}\label{lem:exterior-quotient}
Let $R$ be a commutative ring, $M$ an $R$-module, $M' \subseteq M$ a submodule and $n \ge 0$. Then
\[
\ker\Big(\textstyle\bigwedge^n_R M \twoheadrightarrow \bigwedge^n_R (M/M')\Big) \;=\; \mathrm{span}\big\{ m_1 \wedge \dots \wedge m_n \;:\; m_i \in M' \text{ for some } i \big\}.
\]
\end{lemma}

\begin{proof}
Let $Q$ be the quotient of $\bigwedge^n M$ by the span on the right. Choosing lifts defines an alternating map $(M/M')^n \to Q$, which is well defined because changing a lift changes the wedge by an element with a factor in $M'$. The induced map $\bigwedge^n(M/M') \to Q$ is inverse to the map $Q \to \bigwedge^n(M/M')$.
\end{proof}

\leanref{\lean{exteriorPower_ker_map_mkQ}, \lean{lift_comp_map_mkQ}, \lean{quotientTermMap_ker}, in \lean{DeRham/ExteriorPowerQuotient.lean}.}

\subsection{The algebraic de Rham complex}

The ordinary de Rham complex \cite[Tag~07HY]{stacks} is built following J.~Riou's approach through presentations of modules, ported to the current Mathlib.

\begin{proposition}\label{prop:presentations}
Let $R$ be a commutative ring.
\begin{itemize}
\item[\normalfont{(1)}] The $A$-algebra $B$ is presented by $A[X_b : b \in B]$ with relations $X_{b_1} + X_{b_2} - X_{b_1+b_2}$, $X_{b_1}X_{b_2} - X_{b_1 b_2}$ and $a - X_{a\cdot 1}$. This induces a presentation of the $B$-module $\Omega_{B/A}$ with generators $db$ and relations given by additivity, the Leibniz rule and $d(a) = 0$ for $a \in A$.
\item[\normalfont{(2)}] A presentation of an $R$-module $M$ induces a presentation of $\bigwedge^n_R M$, with generators indexed by $n$-tuples of generators and relations coming from one slot, antisymmetry and alternation.
\item[\normalfont{(3)}] Presentations of finitely many $R$-modules $M_i$ induce a presentation of $\bigotimes_i M_i$.
\end{itemize}
\end{proposition}

\leanref{\lean{Presentation.Tautological.algebraPresentation}, \lean{relation_mem_ker}, \lean{span_relation_le_ker}, \lean{differentialPresentation}, \lean{Presentation.exteriorPower}, \lean{Solution.IsPresentation.exteriorPower}, \lean{Presentation.piTensor}; the files of the last two keep J.~Riou's copyright notice.}

\begin{theorem}\label{thm:derham}
There are $A$-linear maps $d \colon \bigwedge^n_B \Omega_{B/A} \to \bigwedge^{n+1}_B \Omega_{B/A}$, $n \ge 0$, with
\[
d(b_0\, db_1 \wedge \dots \wedge db_n) \;=\; db_0 \wedge db_1 \wedge \dots \wedge db_n
\]
and $d \circ d = 0$. The elements $b_0\,db_1\wedge\dots\wedge db_n$ span $\bigwedge^n_B \Omega_{B/A}$ over $A$.
\end{theorem}

\begin{proof}
Restricting scalars along the presentation of Proposition~\ref{prop:presentations} gives an $A$-linear presentation of $\bigwedge^n_B \Omega_{B/A}$ with generators $b_0\, db_1 \wedge \dots \wedge db_n$. The formula is compatible with each relation; the multiplicative relation uses the Leibniz rule and antisymmetry. On generators $d \circ d = 0$ because $d(1) = 0$.
\end{proof}

\leanref{\lean{DeRhamComplex.presentationDifferentialsDown}, \lean{DeRhamComplex.d}, \lean{DeRhamComplex.d_var}, \lean{DeRhamComplex.d_d}, \lean{deRhamComplex}, \lean{riouSourceGenerator_span}.}

\subsection{The divided power de Rham complex}

The following is \cite[Tag~07HY]{stacks} applied as in \cite[Tag~07HZ]{stacks}.

\begin{theorem}\label{thm:pd-derham}
For every $n \ge 0$ the differential of Theorem~\ref{thm:derham} descends along $\bigwedge^n_B \Omega_{B/A} \twoheadrightarrow \Omega^n_{B,\delta}$. This gives a complex $\Omega^\bullet_{B,\delta}$ of $A$-modules with
\[
d(b\, dg_1 \wedge \dots \wedge dg_n) \;=\; db \wedge dg_1 \wedge \dots \wedge dg_n .
\]
\end{theorem}

\begin{proof}
By Lemma~\ref{lem:exterior-quotient}, the kernel of $\bigwedge^n\Omega_{B/A} \twoheadrightarrow \Omega^n_{B,\delta}$ is spanned over $A$ by the forms
\[
\omega \;=\; b_0\, dg_1 \wedge \dots \wedge \big(dx\dpow{r} - x\dpow{r-1}dx\big) \wedge \dots \wedge dg_n ,
\]
where $x \in J$ and $r \ge 1$, with the relation in one slot. The image of $d\omega$ in $\Omega^{n+1}_{B,\delta}$ is zero: in the quotient $dx\dpow{r} = x\dpow{r-1}dx$, and $dx\dpow{r-1} \wedge dx = x\dpow{r-2}\,dx \wedge dx = 0$, while $d1 = 0$ for $r = 1$.
\end{proof}

\leanref{\lean{pdTermMap_comp_d_pdRelationGenerator_eq_zero}, \lean{quotientTermMap_ker_le_pdRiouSpan}, \lean{pdQuotientRelationStable} (in \lean{DeRham/PDStabilityFull.lean}); \lean{descendLinearMap}, \lean{quotientDifferential}, \lean{quotientDifferential_comp_self}, \lean{deRhamComplexOfQuotient}; \lean{pdDeRhamComplex}, \lean{pdDeRhamComplex_d_generator}.}

\begin{proposition}\label{prop:pd-derham-functorial}
A morphism of divided power algebras $f \colon (B,\delta) \to (D,\delta')$ induces a map of complexes $\Omega^\bullet_{B,\delta} \to \Omega^\bullet_{D,\delta'}$, with $v_1 \wedge \dots \wedge v_n \mapsto f_*v_1 \wedge \dots \wedge f_*v_n$, which is $f$ in degree $0$. This is compatible with identities and composition.
\end{proposition}

\leanref{\lean{pdExteriorMap}, \texttt{pdExteriorMap\_$\iota$Multi}, \lean{pdExteriorMap_smul}, \lean{pdExteriorMap_id}, \lean{pdExteriorMap_comp}, \lean{pdExteriorMap_d_naturality}, \lean{pdDeRhamMap}, \lean{pdDeRhamMap_id}, \lean{pdDeRhamMap_comp}, \lean{zeroEquiv_comp_pdExteriorMap_zero}.}

\subsection{Characteristic zero}

\begin{proposition}\label{prop:char-zero}
Assume that $A$ is a $\Q$-algebra and give $J$ the divided powers $x\dpow{n} \coloneqq x^n/n!$.
\begin{itemize}
\item[\normalfont{(1)}] The map $d \colon B \to \Omega_{B/A}$ is a PD derivation, and $\Omega_{B/A} \simeq \Omega_{B,\delta}$.
\item[\normalfont{(2)}] Every map of $\Q$-algebras sending an ideal into another is a morphism of divided power algebras for these divided powers.
\item[\normalfont{(3)}] The identity of $B$, with the divided powers $x^n/n!$ on $J + IB$, is a compatible divided power envelope of $(B, J)$.
\end{itemize}
\end{proposition}

\leanref{\lean{isPDDerivation_kaehlerDifferential_D}, \lean{ordinaryPDKaehlerEquiv}, in the file \lean{DeRham/CharacteristicZero.lean}; \lean{ratAlgebra_isDPMorphism}, \lean{ratAlgebra_identity_isDPEnvelope}, \lean{ratAlgebra_hasDPEnvelope}, in \lean{Envelope/FreeConstruction/RatAlgebra.lean}.}

%======================================================================
\section{Homological algebra}\label{sec:homological}
%======================================================================

In this section $R$ is a ring and complexes are $\N$-graded cochain complexes of $R$-modules. The statements about bicomplexes are proved elementwise, through the following criteria.

\begin{lemma}\label{lem:elementwise}
Let $K$ and $L$ be complexes and $f \colon K \to L$ a map of complexes.
\begin{itemize}
\item[\normalfont{(1)}] The complex $K$ is exact in degree $n$ if and only if every cocycle of $K^n$ is a coboundary; for $n = 0$ this means $\ker(d^0) = 0$.
\item[\normalfont{(2)}] The map $f$ is a quasi-isomorphism in degree $n$ if and only if for every cocycle $y \in L^n$ there are a cocycle $x \in K^n$ and $z \in L^{n-1}$ with $f(x) = y + dz$, and every cocycle $x \in K^n$ with $f(x)$ a coboundary is a coboundary, with the convention $L^{-1} = 0$.
\end{itemize}
\end{lemma}

\leanref{\lean{exactAt_iff_elementwise}, \lean{quasiIsoAt_iff_elementwise}, \lean{quasiIso_of_elementwise}, in \lean{Homology/Elementwise.lean}.}

\subsection{First-quadrant bicomplexes}

A bicomplex $K$ is a family of $R$-modules $K^{p,q}$, $p, q \in \N$, with commuting differentials $d_h$ and $d_v$. In Lean it is a \lean{RawBi}, whose differentials are indexed like those of Mathlib's \lean{HomologicalComplex}, so that no transport along equalities of indices is needed. Its total complex has terms $\Tot(K)^n \coloneqq \prod_{p+q=n} K^{p,q}$, with differential
\[
(dx)_{p,q} \;=\; d_h x_{p-1,q} + (-1)^p d_v x_{p,q-1},
\]
and transposition with the sign $(-1)^{pq}$ is an isomorphism $\Tot(K^{\mathsf T}) \simeq \Tot(K)$ (\lean{RawBi.totComplex}, \lean{dRaw_dRaw}, \lean{transposeTotIso}, in \lean{Homology/TotM.lean}).

\begin{proposition}\label{prop:acyclic-assembly}
Let $K$ be a first-quadrant bicomplex. If every column of $K$ is exact in every degree, including degree $0$, then $\Tot(K)$ is exact in every degree. The same holds with rows in place of columns.
\end{proposition}

\begin{proof}
By induction on $k$, every cocycle of $\Tot(K)$ is cohomologous to a cocycle vanishing in all columns $p < k$; in degree $n$ take $k = n+1$. In degree $0$ the induction uses the exactness of the columns in degree $0$, and shows that every cocycle vanishes.
\end{proof}

\leanref{\lean{exactAt_totComplex_of_columnsExact}, \lean{exactAt_totComplex_of_rowsExact}, \lean{rawAcyclic_of_columnsExact}, \lean{rawAcyclic_of_rowsExact}.}

\begin{proposition}\label{prop:augmented}
Let $K$ be a first-quadrant bicomplex, $L$ a complex and $\eta_p \colon L^p \to K^{p,0}$ maps commuting with the horizontal differentials, with $d_v \circ \eta = 0$. If for every $p$ the augmented column
\[
0 \longrightarrow L^p \xrightarrow{\ \eta_p\ } K^{p,0} \longrightarrow K^{p,1} \longrightarrow \cdots
\]
is exact, then the map $L \to \Tot(K)$ sending $y \in L^n$ to $\eta_n(y)$ in bidegree $(n,0)$ is a quasi-isomorphism. The same holds for augmented rows $0 \to L^q \to K^{0,q} \to K^{1,q} \to \cdots$.
\end{proposition}

\leanref{\texttt{augRow$\iota$\_quasiIso}, \texttt{augCol$\iota$\_quasiIso}; the second is reduced to the first by transposition.}

\begin{proposition}\label{prop:totmap}
Let $\varphi \colon K \to L$ be a map of first-quadrant bicomplexes that is a quasi-isomorphism on every column. Then $\Tot(\varphi)$ is a quasi-isomorphism.
\end{proposition}

\begin{proof}
The columnwise mapping fibre, with terms $K^{p,q} \oplus L^{p,q-1}$, has exact columns. Its total complex is therefore acyclic by Proposition~\ref{prop:acyclic-assembly}, and it is identified, up to the sign $(-1)^p$ on the second component, with the mapping fibre of $\Tot(\varphi)$, whose acyclicity means that $\Tot(\varphi)$ is a quasi-isomorphism.
\end{proof}

\leanref{\lean{RawBi.totMap}, \lean{totMap_quasiIso}, \lean{Hom.fiber}, \lean{columnsExact_fib}, in \lean{Homology/TotQuasiIso.lean}.}

\begin{proposition}\label{prop:projections}
Let $K$ be a first-quadrant bicomplex. The projection $\Tot(K) \to K^{\bullet,0}$ onto the bottom row is a map of complexes, and it is a quasi-isomorphism if the bicomplex formed by the rows $q \ge 1$ has acyclic total complex. The same holds for the projection $\Tot(K) \to K^{0,\bullet}$ onto the first column and the columns $p \ge 1$.
\end{proposition}

\leanref{\lean{rowZeroProj}, \lean{rowZeroProj_quasiIso} (in \lean{Homology/TotProjection.lean}), \lean{colZeroProj}, \lean{colZeroProj_quasiIso}.}

\subsection{The lemma of Bhatt and de Jong}

Let $K_c$ be a cosimplicial object in $\N$-graded cochain complexes of $R$-modules, and let $K$ be the bicomplex with $K^{p,q} \coloneqq K_c[p]^q$, horizontal differential $\sum_i (-1)^i \delta_i$ and vertical differential the internal one. The following abstracts the double complex argument in the proof of \cite[Theorem~2.12]{bhatt-dejong}, which is also the argument of \cite[Tag~07LG]{stacks}; there it is phrased with the two spectral sequences of the double complex, here it is proved elementwise.

\begin{theorem}\label{thm:bdj}
Let $K_c$ and $K$ be as above.
\begin{itemize}
\item[\normalfont{(1)}] If for every $q \ge 1$ the alternating coface complex of $n \mapsto K_c[n]^q$ is exact, then the projection from $\Tot(K)$ to the alternating coface complex of $n \mapsto K_c[n]^0$ is a quasi-isomorphism.
\item[\normalfont{(2)}] If for every $n$ the map $K_c[0] \to K_c[n]$ induced by the first vertex is a quasi-isomorphism, then the projection $\Tot(K) \to K_c[0]$ is a quasi-isomorphism.
\end{itemize}
\end{theorem}

\begin{proof}
Part (1) follows from Propositions~\ref{prop:projections} and~\ref{prop:acyclic-assembly}. For (2), let $K_c'$ be the constant cosimplicial object on $K_c[0]$ and compare $K_c$ with $K_c'$ through the maps induced by $[n] \twoheadrightarrow [0]$. These are quasi-isomorphisms by two-out-of-three, since composed with the first-vertex maps they give the identity. Proposition~\ref{prop:totmap} applies to the induced map of bicomplexes, and since $[0] \to [0]$ is the identity, the projection $\Tot(K) \to K_c[0]$ factors through it. For $K_c'$ the horizontal differentials are alternately $0$ and the identity, so the columns $p \ge 1$ form a bicomplex with exact rows, and Propositions~\ref{prop:acyclic-assembly} and~\ref{prop:projections} apply.
\end{proof}

\leanref{\lean{RawBi.ofCosimplicial}, \lean{rowZeroComplexIso}, \lean{colZeroComplexIso}, \lean{RawBi.bhattDeJong_row}, \lean{RawBi.bhattDeJong_column}, \lean{toConst}, \lean{toConst_app_zero}, \lean{colZeroProj_eq_totMap_comp}, in \lean{Homology/CosimplicialTot.lean}.}

Only the vertex $0$ is assumed in (2). This is equivalent to the hypothesis on all vertices, since each vertex map $K_c[0] \to K_c[n]$ is a section of the map induced by $[n] \twoheadrightarrow [0]$.

\subsection{Cosimplicial contractions}

\begin{definition}
Let $X$ be a cosimplicial object of a preadditive category. A \emph{cosimplicial contraction} of $X$ is a family of endomorphisms $h_n(\alpha)$ of $X[n]$, for $n \ge 0$ and $\alpha \colon [n] \to [1]$, such that $X(f) \circ h_m(\alpha \circ f) = h_n(\alpha) \circ X(f)$ for every $f \colon [m] \to [n]$, with $h_n(\mathrm{const}_0) = 0$ and $h_n(\mathrm{const}_1) = \mathrm{id}$.
\end{definition}

A cosimplicial contraction is a cosimplicial homotopy from $0$ to the identity in the sense of \cite[Tag~019U]{stacks}; the contraction of the modules of differentials of a Čech nerve is of this form \cite[Example~2.16]{bhatt-dejong}, \cite[Tag~07L9]{stacks}.

\begin{theorem}\label{thm:contraction}
Let $X$ be a cosimplicial object of a preadditive category with a cosimplicial contraction, and let $d$ be its alternating coface differential. There are maps $s_n \colon X[n+1] \to X[n]$ with
\[
s_0 \circ d^0 = \mathrm{id} \;\text{ and }\; d^{n} \circ s_n + s_{n+1} \circ d^{n+1} = \mathrm{id}.
\]
In particular, the alternating coface complex of a cosimplicial $R$-module with a cosimplicial contraction is exact in every degree, including degree $0$.
\end{theorem}

\begin{proof}
Write $H_t \coloneqq h_{n+1}(k \mapsto [t \le k])$ for $0 \le t \le n+2$, an endomorphism of $X[n+1]$, so that $H_0 = \mathrm{id}$ and $H_{n+2} = 0$. The maps $X(\sigma_j) \circ H_{j+1} \colon X[n+1] \to X[n]$, for $0 \le j \le n$, satisfy the identities of a simplicial homotopy from $0$ to the identity on the opposite simplicial object, and Mathlib's \lean{SimplicialObject.Homotopy.toChainHomotopy} gives the chain homotopy. The displayed identities are obtained by unwinding the formulas.
\end{proof}

\leanref{\lean{CosimplicialContraction}, \lean{homotopyComponent}, \lean{CosimplicialContraction.toSimplicialHomotopy}, \lean{contr}, \lean{objD_zero_comp_contr}, \lean{contr_comp_objD_add}, \lean{CosimplicialContraction.exactAt}, in \lean{Homology/CosimplicialContraction.lean}.}

\subsection{Injectives and limits in large functor categories}

The category $\INF(C/A)$ is not small in Mathlib's sense, since its objects live in a larger universe than its morphisms. The injective resolutions and limits needed for derived global sections are therefore constructed directly.

\begin{theorem}\label{thm:enough-injectives}
Let $\mathcal J$ be a category with objects in \lean{Type (v+1)} and morphisms in \lean{Type v}, and let $R$ be a commutative ring that is \lean{Small.{v+1}}.
\begin{itemize}
\item[\normalfont{(1)}] For an object $Z$ of $\mathcal J$ and an injective $R$-module $M$, the diagram $\Coind_Z(M) \colon Y \mapsto \Map(\Hom(Y, Z), M)$ is an injective object of the category of functors $\mathcal J \to \Mod_R$.
\item[\normalfont{(2)}] Every functor $F \colon \mathcal J \to \Mod_R$ embeds into $\prod_Z \Coind_Z(M_Z)$, where $F(Z) \hookrightarrow M_Z$ are chosen embeddings into injective modules. Hence the category of functors $\mathcal J \to \Mod_R$ has enough injectives.
\item[\normalfont{(3)}] The category \lean{ModuleCat.{v+1} R} has limits of shape $\mathcal J$, computed by compatible families, and the limit functor is additive.
\end{itemize}
\end{theorem}

\leanref{\lean{coinduced}, \lean{injective_coinduced}, \lean{coinducedInjectiveEmbedding}, \lean{hasEnoughInjectivesFunctorCategory}, \lean{compatibleFamilyLimitCone}, \lean{hasLimitsOfShapeLargeModuleCat}, \lean{limModuleCatAdditive}, in \lean{Cohomology/FunctorCategoryInjectives.lean}.}

\begin{proposition}\label{prop:precomposition}
Let $\ell \colon \mathcal C \to \mathcal D$ be a functor with a right adjoint $r$, and let $\mathcal E$ be a category with pullbacks. Then precomposition with $\ell$ sends injective objects of the category of functors $\mathcal D \to \mathcal E$ to injective objects of the category of functors $\mathcal C \to \mathcal E$.
\end{proposition}

\begin{proof}
Precomposition with $\ell$ is right adjoint to precomposition with $r$, and the latter preserves monomorphisms because they are detected objectwise.
\end{proof}

\leanref{\lean{precompositionAdjunction}, \lean{injective_whiskeringLeft_obj_of_adjunction}.}

\subsection{Derived functors and derived limits}

\begin{proposition}\label{prop:grothendieck}
Let $\mathcal A$, $\mathcal B$ and $\mathcal C$ be abelian categories with enough injectives, and $f \colon \mathcal A \to \mathcal B$ and $g, g' \colon \mathcal B \to \mathcal C$ additive functors.
\begin{itemize}
\item[\normalfont{(1)}] An isomorphism $g \simeq g'$ induces an isomorphism of right derived functors to the homotopy category.
\item[\normalfont{(2)}] If $f$ is exact and preserves injective objects, then $f$ sends injective resolutions to injective resolutions, and $R(g \circ f)(X) \simeq Rg(f(X))$ for every object $X$ of $\mathcal A$.
\end{itemize}
\end{proposition}

\leanref{\lean{additiveDerivedIso}, \lean{mapInjectiveResolution}, \lean{rightDerivedToHomotopyCategoryCompIso}, in \lean{Ocris/Cohomology.lean}.}

\begin{theorem}\label{thm:derived-left-adjoint}
Let $\ell \colon \mathcal J \to \mathcal P$ be a functor with a right adjoint between categories as in Theorem~\ref{thm:enough-injectives}, and let $G$ be a functor $\mathcal P \to \Mod_R$. Then there is an isomorphism in the homotopy category
\[
R\lim_{\mathcal P}(G) \;\simeq\; R\lim_{\mathcal J}(G \circ \ell).
\]
\end{theorem}

\begin{proof}
The functor $\ell$ is initial, so $\lim_{\mathcal J} (- \circ \ell) \simeq \lim_{\mathcal P}$. Restriction along $\ell$ is exact and preserves injective objects by Proposition~\ref{prop:precomposition}. Proposition~\ref{prop:grothendieck} concludes.
\end{proof}

\leanref{\lean{derivedLimitIsoOfAdjunction}, in \lean{Ocris/Filtration.lean}.}

\spa

Finally, \lean{QuasiIsoEquiv}, in \lean{Ocris/DeRhamComparison.lean}, is the equivalence relation on complexes generated by quasi-isomorphisms, an inductive proposition with constructors \lean{ofHom}, \lean{symm} and \lean{trans}. Equivalent complexes have isomorphic cohomology in every degree (\lean{QuasiIsoEquiv.nonempty_homology_iso}). It replaces an isomorphism in the derived category in the statements below.

%======================================================================
\section{Computing cohomology on \texorpdfstring{$\INF_e$}{INF\_e}}\label{sec:level}
%======================================================================

By Theorem~\ref{thm:adjunction-derived} it suffices to compute derived limits on the single site $\INF_e$ of the Basic Setting. There the argument of \cite[Tag~07JM]{stacks} for a category with a weakly final object and products of pairs applies; we formalised it in the following abstract form.

\subsection{An abstract Čech computation of derived limits}

\begin{definition}\label{def:cech-data}
\emph{Čech nerve data} on a category $\mathcal J$ consist of a cosimplicial object $X$ of $\mathcal J$, nonempty sets $S(Z)$ for the objects $Z$ of $\mathcal J$, and bijections
\[
\Hom_{\mathcal{J}}(X[n], Z) \;\simeq\; S(Z)^{n+1}
\]
natural in $[n]$. For a functor $G \colon \mathcal J \to \Mod_R$, the \emph{Čech complex} $\check C^\bullet(G)$ is the alternating coface complex of $n \mapsto G(X[n])$, and $\check H^n(G) \coloneqq H^n(\check C^\bullet(G))$.
\end{definition}

\leanref{\lean{CechNerveData} with \lean{equiv_map}, and \lean{cech D G}, in \lean{Cohomology/Cech.lean}. Naturality in $Z$ is not required.}

\begin{proposition}\label{prop:cech-degree-zero}
Let $G \colon \mathcal J \to \Mod_R$ be a functor. The augmentation $\lim(G) \hookrightarrow G(X[0])$ is injective with image $\ker(d^0)$. In particular $\lim(G) \simeq \check H^0(G)$.
\end{proposition}

\begin{proof}
Given two maps $g, g' \colon X[0] \to Z$, the bijections provide $h \colon X[1] \to Z$ whose two faces are $g$ and $g'$. Hence a cocycle in $G(X[0])$ has the same image under all maps out of $X[0]$, and defines a compatible family. The augmentation is injective because every object $Z$ receives a map from $X[0]$, the set $S(Z)$ being nonempty.
\end{proof}

\leanref{\lean{CechNerveData.exists_pair}, \lean{augment_injective}, \lean{exists_limit_of_cocycle}, \lean{augment_comp_d}, \lean{cocycle_map_eq}, \lean{limitIso} (limits as compatible families).}

\begin{proposition}\label{prop:cech-injective}
If $G$ is an injective object of the category of functors $\mathcal J \to \Mod_R$, then $\check C^\bullet(G)$ is exact in every positive degree.
\end{proposition}

\begin{proof}
Consider the functor
\[
\mathrm{big}(G) \colon Y \longmapsto \prod_{Z} \Map\big(\Hom_{\mathcal{J}}(Y, Z),\, G(Z)\big)
\]
with the monomorphism $G \hookrightarrow \mathrm{big}(G)$ sending $a \in G(Y)$ to $(g \mapsto G(g)(a))_Z$. The Čech complex of $\mathrm{big}(G)$ has an explicit contraction in positive degrees, obtained by prepending a chosen element of $S(Z)$. Since $G$ is injective the monomorphism has a retraction, so exactness transfers to $G$.
\end{proof}

\leanref{\lean{bigCoinduced}, \lean{toBig}, \lean{mono_toBig}, \lean{bigContraction}, \lean{big_cocycle_eq}, \lean{bigRetraction}, \lean{exists_cech_d_eq_of_injective}.}

\begin{theorem}\label{thm:cech-zigzag}
Let $\mathcal J$ be equipped with Čech nerve data, with objects in \lean{Type (v+1)} and morphisms in \lean{Type v}, let $R$ be a commutative ring, let $F \colon \mathcal J \to \Mod_R$ be a functor and $F \to I^\bullet$ an injective resolution. There is a zigzag of quasi-isomorphisms
\[
\lim(I^\bullet) \longrightarrow \Tot\big(I^q(X[n])\big) \longleftarrow \check{C}^\bullet(F).
\]
\end{theorem}

\begin{proof}
The bicomplex $I^q(X[n])$ has two augmentations. The augmented rows $\lim(I^q) \to I^q(X[0]) \to I^q(X[1]) \to \cdots$ are exact by Propositions~\ref{prop:cech-degree-zero} and~\ref{prop:cech-injective}. The augmented columns $F(X[n]) \to I^0(X[n]) \to I^1(X[n]) \to \cdots$ are exact because a resolution in a functor category is exact objectwise. Proposition~\ref{prop:augmented} makes both augmentations quasi-isomorphisms.
\end{proof}

\leanref{\lean{cechBicomplex}, \lean{limAug}, \lean{limAug_rowsExact}, \lean{resAug}, \lean{resAug_columnsExact}, \lean{res_exact_zero}, \lean{res_exact_succ}, \lean{res_mono}, \lean{cechZigzag}.}

Combined with Mathlib's computation of right derived functors by injective resolutions, this gives the following statement.

\begin{theorem}\label{thm:cech-derived}
Let $\mathcal J$ be a category with objects in \lean{Type (v+1)} and morphisms in \lean{Type v}, equipped with Čech nerve data, and let $R$ be a commutative ring that is \lean{Small.{v+1}}. For every functor $F \colon \mathcal J \to \Mod_R$ and every $n \ge 0$,
\[
R^n \lim(F) \;\simeq\; \check H^n(F).
\]
In particular $\lim(F) \simeq \check H^0(F)$, and injective functors have vanishing Čech cohomology in positive degrees.
\end{theorem}

\leanref{\lean{Cohomology.Cech.rightDerivedLimIsoCech}, \lean{limIsoCechHomologyZero}, \lean{cech_homology_isZero_of_injective}, \lean{Zigzag.homologyIso}, \lean{limComplexHomologyIsoCech}, in \lean{Cohomology/CechDerived.lean}.}

\subsection{The Čech--Alexander complex}

On $\INF_e$, the Čech nerve data consist of the cosimplicial object $X_e(\bullet)$, the sets of sections $S(Z)$, and the bijections of Proposition~\ref{prop:cech-hom} (\lean{Level.cechNerveData}). Applied to $\Ocris$ they give the Čech complex $D(\bullet)$ (\lean{ocrisCechComplex}) and the zigzag of Theorem~\ref{thm:cech-zigzag} (\lean{ocrisCechZigzag}).

\begin{theorem}\label{thm:cech-alexander}
Assume the Basic Setting and let $E$ be an envelope family. For every $n \ge 0$ there are isomorphisms of $A$-modules
\[
H^n\big(R\Gamma(\INF(C/A), \Ocris)\big) \;\simeq\; R^n\lim{}_{\INF_e}(\Ocris) \;\simeq\; H^n\big(D(0) \to D(1) \to D(2) \to \cdots\big).
\]
In particular the cohomology of the complex $D(\bullet)$ does not depend on the choice of $e$ with $(p^e)! = 0$ in $A$.
\end{theorem}

\begin{proof}
The first isomorphism is Theorem~\ref{thm:adjunction-derived}, and the second is Theorem~\ref{thm:cech-derived} for the Čech nerve data on $\INF_e$.
\end{proof}

\leanref{\lean{EnvelopeFamily.crystallineCohomologyIsoCechAlexander}, \lean{levelRightDerivedIsoCechAlexander}, \lean{cechAlexander_homology_independent}, in \lean{Ocris/CechAlexander.lean}; \lean{InfinitesimalSite.crystallineCohomologyIsoCechAlexander}, for \lean{envelopeFamily}, in \lean{MainTheorems.lean}; also \lean{ocrisDerivedGlobalSectionsIsoLimComplex}, \lean{levelOcrisDerivedGlobalSectionsIsoLimComplex}. The statements in \lean{Ocris/CechAlexander.lean} are proved on \lean{Level A C m} for every $m > 0$ with $m! = 0$ in $A$, and \lean{cechAlexander_homology_independent} compares two such exponents.}

The terms $D(n)$ are envelopes of the polynomial presentations $P_n \twoheadrightarrow C$ by Proposition~\ref{prop:cech-envelope}. Theorem~\ref{thm:cech-alexander} needs neither the Poincaré lemma nor the vanishing lemma of the next section.

%======================================================================
\section{The de Rham comparison}\label{sec:derham}
%======================================================================

We keep the notation of the previous section: we assume the Basic Setting, and $D(n) = \Ocris(X_e(n))$.

\subsection{The bicomplex}

The assignment $n \mapsto \Omega^\bullet_{D(n)}$ is a cosimplicial object in cochain complexes of $A$-modules (\lean{levelCechDeRham}). Let $K$ be the associated bicomplex, with $K^{n,q} \coloneqq \Omega^q_{D(n)}$, horizontal differential $d_h = \sum_{i=0}^{n+1} (-1)^i \, \delta_i^*$ and vertical differential $d_v = d$, in Lean \lean{RawBi.ofCosimplicial}. The comparison runs through the chain
\begin{equation}\label{eq:chain}
\lim(I^\bullet) \longrightarrow \Tot\big(I^q(X_e(n))\big) \longleftarrow \check{C}^\bullet(\Ocris) \;\simeq\; \Omega^0_{D(\bullet)} \;\longleftarrow\; \Tot(K) \;\longrightarrow\; \Omega^\bullet_{D(0)} ,
\end{equation}
where $\Ocris|_{\INF_e} \to I^\bullet$ is an injective resolution in the category of functors on $\INF_e$. The first two arrows are Theorem~\ref{thm:cech-zigzag}, the isomorphism identifies $0$-forms with functions (\lean{levelDeRhamZeroIso}, \lean{rowZeroComplexIso}), and the last two arrows are the projections of Theorem~\ref{thm:bdj}. The theorem \lean{affineDeRhamComparison_of_quasiIso_of_exact} assembles \eqref{eq:chain} from the hypotheses of the two parts of Theorem~\ref{thm:bdj}, called \lean{PositiveCechExact} and \lean{FirstVertexQuasiIso}. The rest of the section proves both.

\subsection{Presented envelopes}

To construct operators on $\Omega_{D(n)}$ we use only the universal property.

\begin{definition}
Let $\sigma$ be a set and $J \subseteq A[x_s : s \in \sigma]$ an ideal. A \emph{presented envelope} of $J$ is an $A$-algebra map $\iota \colon A[x_s : s \in \sigma] \to D$ together with divided powers $\delta$ on an ideal $K \supseteq JD + ID$, compatible with $\gamma$, such that $\iota$ is initial among maps to compatible divided power algebras sending $J$ into the divided power ideal. Put $dx_s \coloneqq d(\iota(x_s))$.
\end{definition}

\begin{proposition}\label{prop:presented}
Let $(\iota, D, \delta)$ be a presented envelope of $J$.
\begin{itemize}
\item[\normalfont{(1)}] Morphisms of divided power algebras and PD derivations out of $D$ are determined by their values on the elements $\iota(x_s)$.
\item[\normalfont{(2)}] The $D$-module $\Omega_{D,\delta}$ is generated by the $dx_s$. Hence $D$-linear maps out of $\Omega_{D,\delta}$ are determined by their values on the $dx_s$.
\item[\normalfont{(3)}] An $A$-subalgebra $D' \subseteq D$ containing every $\iota(x_s)$ and such that $\delta_n(y) \in D'$ for all $y \in D' \cap K$ and $n \ge 1$ is equal to $D$.
\item[\normalfont{(4)}] An $A$-algebra map $\varphi \colon A[x_s] \to A[x_{s'}]$ with $\varphi(J) \subseteq J'$ induces a morphism of divided power algebras between presented envelopes of $J$ and $J'$.
\item[\normalfont{(5)}] For every family $(m_s)$ of elements of a $D$-module there is a unique PD derivation $\theta$ of $D$ with $\theta(\iota(x_s)) = m_s$.
\end{itemize}
\end{proposition}

\begin{proof}
Part (1) follows from initiality and Theorem~\ref{thm:derivation-extension}. For (2), the composite of $d$ with the projection to $\Omega_{D,\delta}/\langle dx_s\rangle$ is a PD derivation vanishing on the $\iota(x_s)$, hence zero by (1). For (3), the divided powers restrict to $D'$, initiality gives $\psi \colon D \to D'$, and the inclusion composed with $\psi$ is the identity by (1). Parts (4) and (5) follow from initiality and Theorem~\ref{thm:derivation-extension}.
\end{proof}

\leanref{\lean{IsPresentedEnvelope}, \lean{PresentedEnvelope.hom_ext}, \lean{PresentedEnvelope.pdDerivation_ext}, \lean{span_dx}, \lean{hom_ext_dx}, \lean{PresentedEnvelope.eq_top}, \lean{DividedPowers.restrictSubalgebra}, \lean{PresentedEnvelope.lift}, \lean{lift_X}, \lean{PresentedEnvelope.extendDer}, \lean{extendDer_spec}, in \lean{Ocris/PresentedEnvelope.lean} and \lean{Ocris/PresentedEnvelopeTools.lean}; \lean{cech_isPresentedEnvelope} in \lean{Ocris/CechVanishing.lean}.}

\subsection{The vanishing lemma}

The following is \cite[Tag~07LA]{stacks}; its proof here uses no freeness of $\Omega_{D(n)}$.

\begin{proposition}\label{prop:vanishing}
For every $q \ge 1$, the alternating coface complex of $n \mapsto \Omega^q_{D(n)}$ is exact.
\end{proposition}

\begin{proof}
By Proposition~\ref{prop:presented}(2) and~(5), for every predicate $S$ on a set of variables there is a unique $D$-linear map $\partial_S \colon \Omega_{D,\delta} \to \Omega_{D,\delta}$ with
\[
\partial_S(dx_s) \;=\; \begin{cases} dx_s & \text{if } S(s), \\ 0 & \text{otherwise.} \end{cases}
\]
If a morphism of presented envelopes $g$ sends $\iota(x_s)$ to $\iota'(x_{r(s)})$, then $g_* \circ \partial_{S' \circ r} = \partial'_{S'} \circ g_*$, and the same holds for exterior powers. For $\alpha \colon [n] \to [1]$ let $S_\alpha(i,c)$ be the predicate $\alpha(i) = 1$ on the variables of $P_n$, and put
\[
h_n(\alpha) \;\coloneqq\; \textstyle\bigwedge^{q} \partial_{S_\alpha} \colon \Omega^{q}_{D(n)} \longrightarrow \Omega^{q}_{D(n)} .
\]
This is natural in $([n], \alpha)$, $h_n(\mathrm{const}_1) = \mathrm{id}$, and $h_n(\mathrm{const}_0) = 0$ because $q \ge 1$. Theorem~\ref{thm:contraction} applies.
\end{proof}

\leanref{\lean{partialOp}, \lean{partialOp_dx}, \lean{partialOp_of_forall}, \lean{partialOp_of_forall_not}, \lean{partialExt}, \lean{partialExt_natural}, \lean{cech_comparison_X}, \lean{cechContraction}, \lean{positiveCechExact}, in \lean{Ocris/PresentedEnvelope.lean} and \lean{Ocris/CechVanishing.lean}.}

\subsection{The one-variable Poincaré lemma}

\begin{definition}
Let $(B, J, \delta)$ be a divided power algebra. A \emph{one-variable structure} on $B$ consists of an idempotent morphism of divided power algebras $\varepsilon \colon B \to B$, an element $\xi \in J$ with $\varepsilon(\xi) = 0$, and a PD derivation $\partial \colon B \to B$ with $\partial(\xi) = 1$ and $\partial \circ \varepsilon = 0$, such that every $y \in B$ is a finite sum $\sum_k \xi\dpow{k}\, \varepsilon(c_k)$ with $c_k \in B$.
\end{definition}

This models $B = B_0\langle \xi \rangle$ with $B_0 \coloneqq \varepsilon(B)$, without asserting that the divided powers of $\xi$ form a basis. The following generalises the one-variable case of \cite[Tag~07LD]{stacks}; its proof follows \cite[Tags~07Q7, 07Q8]{stacks}.

\begin{theorem}\label{thm:one-variable}
Let $(\varepsilon, \xi, \partial)$ be a one-variable structure on $(B, J, \delta)$ and let $\varepsilon^*$ be the induced endomorphism of $\Omega^\bullet_{B,\delta}$. If $\omega \in \Omega^0_{B,\delta}$ is a cocycle, then $\omega = \varepsilon^*\omega$. If $\omega \in \Omega^{m+1}_{B,\delta}$ is a cocycle, then $\omega - \varepsilon^*\omega = d\eta$ for some $\eta \in \Omega^m_{B,\delta}$.
\end{theorem}

\begin{proof}
For a $1$-form $v$ write $v \wedge -$ for the wedge product, and for a $B$-linear form $\phi$ on $\Omega_{B,\delta}$ let $\iota_\phi$ be the interior product, defined through Mathlib's \lean{CliffordAlgebra.contractLeft} for the zero quadratic form. They satisfy $v \wedge v \wedge x = 0$ and
\[
\iota_\phi(v \wedge x) \;=\; \phi(v)\,x - v \wedge \iota_\phi x,
\]
as well as the Leibniz rule $d(fx) = df \wedge x + f\,dx$.
Let $\iota_\partial$ be the interior product with the form induced by $\partial$ (Proposition~\ref{prop:pd-kaehler}). One checks that $\xi\dpow{a}\,\xi\dpow{b} = \tbinom{a+b}{a}\,\xi\dpow{a+b}$, that $\varepsilon(\xi\dpow{k}) = \delta_{k0}$ and that $d\xi\dpow{k+1} = \xi\dpow{k}\, d\xi$, and moreover that $(\varepsilon^*)^2 = \varepsilon^*$, $d\varepsilon^* = \varepsilon^* d$ and $\iota_\partial \circ \varepsilon^* = 0$.

\spa

By the Taylor property, write $\omega = \sum_k \xi\dpow{k} (\alpha_k + d\xi \wedge \beta_k)$ with $\alpha_k$ and $\beta_k$ in the image of $\varepsilon^*$. Replacing $\omega$ by $\omega - d\big(\sum_k \xi\dpow{k+1} \beta_k\big)$, one may assume $\omega = \sum_k \xi\dpow{k} \rho_k$ with $\rho_k$ in the image of $\varepsilon^*$. The coefficients are recovered as
\[
\rho_a \;=\; \varepsilon^*\Big( (\iota_\partial \circ d)^a\, \omega \Big),
\]
and since $d\omega = 0$ every $\rho_a$ with $a \ge 1$ vanishes. Hence the modified form is equal to $\rho_0$, which is fixed by $\varepsilon^*$. Since $\varepsilon^*$ commutes with $d$, the original form satisfies $\omega - \varepsilon^*\omega = d\eta$ with $\eta \coloneqq \sum_k \xi\dpow{k+1} \beta_k$, since $\varepsilon^*\eta = 0$; in degree $0$ there is no $\beta_k$ and $\omega = \varepsilon^*\omega$.
\end{proof}

\leanref{\lean{OneVariable}, \lean{OneVariable.cocycle_zero}, \lean{OneVariable.cocycle_succ}, \lean{OneVariable.isBd_sub}, \lean{exists_repr}, \lean{extract}, \lean{cocycle_pure}, \lean{OneVariable.dd_dd} (in \lean{DeRham/OneVariablePoincare.lean}); \lean{wdg}, \lean{wdg_wdg_self}, \lean{ctr}, \lean{ctr_wdg_succ}, \lean{dd_smul} (in \lean{DeRham/FormCalculus.lean}).}

\subsection{The relative Poincaré lemma}

The following is the case of \cite[Tag~07LD]{stacks} needed here.

\begin{proposition}\label{prop:poincare}
For every $n \ge 0$, the map $\Omega^\bullet_{D(0)} \to \Omega^\bullet_{D(n)}$ induced by the first vertex is a quasi-isomorphism.
\end{proposition}

\begin{proof}
The retraction $[n] \twoheadrightarrow [0]$ gives a left inverse, so the map is injective on cohomology. For surjectivity it suffices to show that every cocycle $\omega$ of $\Omega^\bullet_{D(n)}$ is cohomologous to $\varepsilon_{\mathrm{all}}^* \omega$, where $\varepsilon_{\mathrm{all}}$ is the endomorphism of $D(n)$ induced by $x_{i,c} \mapsto x_{0,c}$.

\spa

Let $t = (i,c)$ be a variable with $i \ne 0$. Let $\varepsilon_t$ be the endomorphism of $D(n)$ induced by $x_t \mapsto x_{0,c}$ and fixing the other variables (Proposition~\ref{prop:presented}(4)), put $\xi_t \coloneqq x_t - x_{0,c}$, and let $\partial_t$ be the PD derivation with $\partial_t(x_s) = [s = t]$ (Proposition~\ref{prop:presented}(5)). The identities $\partial_t \circ \varepsilon_t = 0$ and $\varepsilon_t^2 = \varepsilon_t$ follow from Proposition~\ref{prop:presented}(1). The set of finite sums $\sum_k \xi_t\dpow{k}\varepsilon_t(c_k)$ is an $A$-subalgebra containing all variables and stable under divided powers, so by Proposition~\ref{prop:presented}(3) it is $D(n)$. Thus $(\varepsilon_t, \xi_t, \partial_t)$ is a one-variable structure.

\spa

For a set $S$ of variables let $\varepsilon_S$ be the endomorphism of $D(n)$ induced by $x_{i,c} \mapsto x_{0,c}$ for $(i,c) \in S$ and fixing the other variables; thus $\varepsilon_{\mathrm{all}} = \varepsilon_S$ for $S$ the set of all variables, and $\varepsilon_{S \cup S'} = \varepsilon_S \circ \varepsilon_{S'}$. The variables are indexed by $[n] \times C$, which is infinite when $C$ is. Every form $\omega$ has a finite set $F$ of variables such that $\varepsilon_{S}^*\omega = \omega$ whenever $S$ is disjoint from a finite set containing $F$; this is proved by induction on $\omega$, using Proposition~\ref{prop:presented}(3) for the elements of $D(n)$. Hence $\varepsilon_{\mathrm{all}}^*\omega = \varepsilon_F^*\omega$. By induction on $F$, composing the homotopies of Theorem~\ref{thm:one-variable} for the variables $t \in F$, the form $\omega - \varepsilon_F^*\omega$ is a coboundary when $\omega$ is a cocycle.
\end{proof}

\leanref{\lean{oneVariable}, \lean{taylorSubalgebra_εS_eq_top}, \lean{exists_support}, \lean{exists_support_form}, \lean{isBd_sub_finset}, \lean{isBd_sub_all}, \lean{comparison_const_eq_εS}, \lean{firstVertexQuasiIso}, in \lean{Ocris/CechPoincare.lean}; \lean{taylorSubalgebra}, \lean{taylorSubalgebra_dpow_mem} in \lean{DeRham/OneVariablePoincare.lean}.}

\begin{remark}
Neither Proposition~\ref{prop:vanishing} nor Proposition~\ref{prop:poincare} uses a divided power polynomial algebra, nor freeness of $D(n)$ over $D(0)$ or of $\Omega_{D(n)}$. The only access to $D(n)$ is its universal property, through Proposition~\ref{prop:presented}.
\end{remark}

\subsection{The main theorem}

\begin{theorem}\label{thm:main}
Assume the Basic Setting and let $E$ be an envelope family. Let $\Ocris|_{\INF_e} \to I^\bullet$ be an injective resolution in the category of functors on $\INF_e$. The complex $\lim(I^\bullet)$, which computes $R\Gamma(\INF(C/A), \Ocris)$ by Theorem~\ref{thm:adjunction-derived}, is connected to $\Omega^\bullet_{D(0)}$ by a zigzag of quasi-isomorphisms, and for every $k \ge 0$ there is an isomorphism of $A$-modules
\[
H^k\big(R\Gamma(\INF(C/A), \Ocris)\big) \;\simeq\; H^k\big(\Omega^\bullet_{D(0)}\big).
\]
\end{theorem}

\begin{proof}
Apply \lean{affineDeRhamComparison_of_quasiIso_of_exact} to the chain \eqref{eq:chain}, with Proposition~\ref{prop:vanishing} for Theorem~\ref{thm:bdj}(1) and Proposition~\ref{prop:poincare} for Theorem~\ref{thm:bdj}(2), and then Theorem~\ref{thm:adjunction-derived}.
\end{proof}

By Proposition~\ref{prop:cech-envelope}, $D(0)$ is a compatible divided power envelope of $A[x_c : c \in C] \twoheadrightarrow C$, so the right-hand side does not depend, up to isomorphism, on the choice of $e$ with $(p^e)! = 0$ in $A$. In Lean the theorem is \lean{EnvelopeFamily.crystallineCohomologyIsoDeRham}, in \lean{Ocris/CechPoincare.lean}, for an arbitrary envelope family \lean{E} and an arbitrary exponent \lean{n} with $n! = 0$ in $A$:

\begin{lstlisting}
variable (E : EnvelopeFamily (A := A) (C := C) ($\gamma$ := $\gamma$))
variable (p : $\mathbb{N}$) [Fact p.Prime] [IsCrystallineBase p A I]
  (hIC : I.map (algebraMap A C) = $\bot$)
variable (n : $\mathbb{N}$) [NeZero n] (hn : ((Nat.factorial n : $\mathbb{N}$) : A) = 0)

include p hIC hn in
theorem crystallineCohomologyIsoDeRham (k : $\mathbb{N}$) :
    Nonempty (((HomotopyCategory.homologyFunctor (ModuleCat.{u + 1} A)
        (ComplexShape.up $\mathbb{N}$) k).obj
        (ocrisDerivedGlobalSections E)) $\cong$
      ((ocrisDeRhamComplex E).obj (Level.cechObj n 0).obj).homology k)
\end{lstlisting}

Here \lean{n} is the exponent of the thickenings, not a Čech degree, and Theorem~\ref{thm:main} is the case $n = p^e$, with \lean{Level.cechObj (p ^ e) 0} the object $X_e(0)$. The complex-level form is \lean{affineDeRhamComparison}, which asserts \lean{QuasiIsoEquiv} between \lean{Cohomology.Cech.limComplex} of the injective resolution of \lean{E.levelOcrisModule p hIC n} and $\Omega^\bullet_{D(0)}$. Combining Theorems~\ref{thm:main} and~\ref{thm:cech-alexander}, the complexes $D(\bullet)$ and $\Omega^\bullet_{D(0)}$ have isomorphic cohomology. The file \lean{MainTheorems.lean} states the three theorems, under the same names in the namespace \lean{InfinitesimalSite}, for the envelope family \lean{envelopeFamily} given by Theorem~\ref{thm:envelope-existence}, so that they take no envelope family as input, on the Frobenius level \lean{Level A C (p ^ e)} and under the hypothesis \texttt{he : ((Nat.factorial (p \textasciicircum{} e) : $\mathbb{N}$) : A) = 0}.

%======================================================================
\section{How the formalisation was produced}\label{sec:produced}
%======================================================================

We used Claude, an AI system developed by Anthropic, to autoformalise the proof, following the Stacks Project \cite{stacks}, chapter \emph{Crystalline Cohomology}, and the double complex lemma of Bhatt and de Jong \cite{bhatt-dejong}. The mathematical design is due to the author: the use of the site of infinitesimal thickenings whose kernel has a uniform nilpotence exponent, in place of a site whose objects carry divided powers, and the strategy of the proof. The Lean code of the library, except two files on presentations of modules due to J.~Riou (Table~\ref{tab:infra}), was written by Claude and checked by the Lean compiler; it builds on Mathlib \cite{mathlib}. The text of this article was also drafted with Claude, from the formalisation and under the author's direction.

\spa

The formalisation follows four principles. The site is presentation-free: polynomial presentations appear in theorems about the site, never in the definition of the site or of $\Ocris$. Divided power envelopes are used only through their universal property \lean{IsDPEnvelope}. No data is left unproved: a \lean{Type} or an instance that statements quantify over is never defined by \texttt{sorry}. Finally, $p$ is nilpotent and no $p$-adic completion appears. In particular, the proofs of Propositions~\ref{prop:vanishing} and~\ref{prop:poincare} use only the universal property of envelopes, and no freeness statement.

\subsection{Components of the library}

Table~\ref{tab:infra} lists the main components and their sources.

\begin{table}[ht]
\centering
\small
\renewcommand{\arraystretch}{1.25}
\begin{tabularx}{\linewidth}{@{}>{\raggedright\arraybackslash}X >{\raggedright\arraybackslash}p{0.33\linewidth} >{\raggedright\arraybackslash}p{0.2\linewidth}@{}}
\toprule
Component & Where & Source \\
\midrule
Divided powers and their morphisms, \lean{DividedPowers}, \lean{IsDPMorphism} & \lean{Mathlib.RingTheory.DividedPowers} & Mathlib, by Chambert-Loir and de Frutos-Fernández \cite{cl-df} \\
Divided power algebra, \lean{DividedPowerAlgebra} & \lean{Mathlib.RingTheory.DividedPowerAlgebra} & Mathlib, by Chambert-Loir and de Frutos-Fernández \cite{cl-df-laws} \\
Exterior algebra, \lean{CliffordAlgebra.contractLeft}, alternating coface complexes, \lean{SimplicialObject.Homotopy.toChainHomotopy}, derived functors & Mathlib & Mathlib \cite{mathlib} \\
Presentations of exterior powers and tensor products & \lean{DeRham/Presentation/ExteriorPowerExplicit.lean}, \lean{DeRham/Presentation/PiTensor.lean} & J.~Riou \\
The infinitesimal sites $\INF(C/A)$ and $\INF_e$, the truncation adjunction, coproducts, Čech nerves & \lean{InfinitesimalSite/} & This formalisation \\
Envelopes: universal property, choices, functoriality, invariance under truncation & \lean{PDDifferentials.lean}, \lean{InfinitesimalSite/Envelope.lean}, \lean{Envelope/}, \lean{Ocris/Filtration.lean} & This formalisation \\
Existence of divided power envelopes & \lean{Envelope/Construction.lean} & This formalisation, following \cite[Tag~07H8]{stacks} \\
Divided power algebras and Roby monomials, towards an explicit construction of envelopes & \lean{Envelope/FreeConstruction/}, \lean{InfinitesimalSite/EnvelopeConstruction.lean} & This formalisation \\
Divided power differentials and de Rham complexes & \lean{PDDifferentials.lean}, \lean{PDDifferentials/}, \lean{DeRham/} & This formalisation, following Riou's approach \\
$\Ocris$, derived global sections, Čech--Alexander and de Rham comparisons & \lean{Ocris.lean}, \lean{Ocris/} & This formalisation \\
Bicomplexes, Bhatt--de Jong lemma, cosimplicial contractions, Čech computation of derived limits, injectives in functor categories & \lean{Homology/}, \lean{Cohomology/} & This formalisation \\
\bottomrule
\end{tabularx}
\caption{Components of the library and their sources.}
\label{tab:infra}
\end{table}

\subsection{Verification}

The library consists of 83 Lean files and about 23000 lines. It builds with \lean{lake build} and contains no \texttt{sorry}; it follows the Mathlib naming conventions and builds without warnings under Mathlib's standard set of style linters. The command \lean{#print axioms} for \lean{crystallineCohomologyIsoDeRham}, \lean{affineDeRhamComparison} and \lean{crystallineCohomologyIsoCechAlexander}, both for an arbitrary envelope family and for \lean{envelopeFamily}, and for \lean{hasDPEnvelope}, lists only \lean{propext}, \lean{Classical.choice} and \lean{Quot.sound}. The continuous integration of the repository checks the build, the absence of \texttt{sorry} and these axioms on every change, with \lean{scripts/CheckAxioms.lean}.

%======================================================================
\section{Scope and what remains}\label{sec:scope}
%======================================================================

Theorem~\ref{thm:main} is an affine statement with trivial coefficients. In this section we explain what its hypotheses restrict and what remains to be done.

\subsection{What the hypotheses mean}

Since $p$ is nilpotent in $A$, no $p$-adic completion and no limit over powers of $p$ appears. The hypothesis $(p^e)! = 0$ in $A$ is a quantitative form of this nilpotence: by Remark~\ref{rem:valuation} it is equivalent to $p^{(p^e-1)/(p-1)} = 0$ in $A$, and it holds for $e = N$ when $p^N = 0$. It enters the proofs of the comparison theorems only through Lemma~\ref{lem:pd-exponent}, in Lemma~\ref{lem:truncation-envelope}, so that truncation does not change divided power envelopes. The comparison theorems are proved in Lean on the thickenings of exponent $m$, for every $m > 0$ with $m! = 0$ in $A$, and the fields of \lean{IsCrystallineBase}, namely $p \in I$ and the nilpotence of $p$, are not used in their proofs. The condition $I \cdot C = 0$ means that $\gamma$ extends to $C$ trivially, on the zero ideal, so it is a special case of the usual compatibility condition.

\spa

The base is $\Spec(C)$, the objects are affine thickenings $B \twoheadrightarrow C$, the topology is chaotic, and cohomology is the right derived functor of the limit. Only the structure sheaf $\Ocris$ is treated: no crystal of modules and no de Rham complex with connection appears.

\spa

The theorems of \lean{Ocris/} take an envelope family $E$, used only through \lean{IsDPEnvelope}; those of \lean{MainTheorems.lean} use the family \lean{envelopeFamily} of Theorem~\ref{thm:envelope-existence}, and the cohomology does not depend on $E$. The comparison is stated on $\INF_e$ for some $e$ with $(p^e)! = 0$ in $A$; by Proposition~\ref{prop:cech-envelope}, $D(0)$ is an envelope of $A[x_c : c \in C] \twoheadrightarrow C$ itself, so the right-hand side does not depend on $e$.

\subsection{An explicit description of envelopes}\label{subsec:existence}

The existence proof of Theorem~\ref{thm:envelope-existence} realises the envelope inside a product indexed by all admissible pairs, and gives no explicit description of it. Such a description is not needed for the comparison theorems, but it would be needed to compute: for instance, to show that the envelope of a polynomial presentation in which $J$ is generated by variables is a divided power polynomial algebra.

\spa

The classical description is as a quotient of a divided power algebra, $D \simeq \Gamma_B(J)/J_{12}$ \cite{berthelot-ogus,roby65}; see also \cite{cl-df,cl-df-laws} for the formalisation of divided power algebras in Lean. The folder \lean{Envelope/FreeConstruction/} contains results in this direction. For the divided power algebra $\Gamma_R(M)$, the library proves that the positive ideal $\Gamma^+$ is the kernel of the augmentation, that $\Gamma/\Gamma^+ \simeq R$, and that divided powers on $\Gamma^+$ are determined by their values on the generators $\gamma_k(m)$ (\lean{positiveIdeal_eq_augmentation_ker}, \lean{positiveQuotientEquiv}, \lean{dividedPowers_eq_of_positiveGenerator_agreement}). When $M$ is free with a basis $b$, it proves that the Roby monomials \cite{roby63} span each graded piece and are linearly independent in degrees $0$ and $1$ (\lean{robyGradeMonomial_span_eq_top}, \lean{robyGradeZeroBasis}, \lean{robyGradeMonomial_degree_one_linearIndependent}). Assuming that the Roby monomials form a basis and that Roby's formula satisfies the axioms of divided powers, it constructs divided powers on $\Gamma^+$ and proves that they are unique (\lean{dividedPowers_unique}). Assuming divided powers on $\Gamma^+$ satisfying Roby's formula, it proves the binomial expansion of the divided powers of the second family of relations (\lean{rel2_dpow_expansion}). If moreover $J_{12} \cap \Gamma^+$ is a sub-divided power ideal, the quotient $\Gamma_B(J')/J_{12}$, with $J' \coloneqq J + IB$, inherits divided powers (\lean{IncludedEnvelopeConstructionData.quotientDividedPowers}, \lean{IncludedEnvelopeConstructionData.quotientMap_isDPMorphism}). Finally, if for every target the map out of $\Gamma_B(J')$ given by its universal property is a morphism of divided power algebras vanishing on $J_{12}$, the induced map is the unique compatible one (\lean{QuotientUniversalProperty.compatible_map_eq_quotientMap}).

\spa

The following is not proved: the linear independence of Roby monomials in degrees at least $2$; Roby's theorem that $\Gamma^+$ carries divided powers, whose formalisation through polynomial laws is in progress \cite{cl-df-laws}; the stability of $J_{12} \cap \Gamma^+$ under these divided powers; and the compatibility of the maps out of $\Gamma_B(J')$ with divided powers. Hence the envelope of Theorem~\ref{thm:envelope-existence} is not identified with $\Gamma_B(J')/J_{12}$.

\subsection{Comparison with the crystalline site}

The library computes crystalline cohomology as $R\Gamma(\INF(C/A), \Ocris)$ and does not define Berthelot's crystalline site of $\Spec(C)$ over $\Spec(A)$, whose objects are divided power thickenings $U \hookrightarrow T$ with $U \subseteq \Spec(C)$ open, with a Zariski or étale topology. That its cohomology agrees with $R\Gamma(\INF(C/A), \Ocris)$ is classical and is not formalised. It follows from Theorem~\ref{thm:main} and the classical comparison with $\Omega^\bullet_D$ \cite[Tag~07LG]{stacks}, or from the description of the cohomology of the crystalline site of an affine scheme as a derived limit over affine divided power thickenings \cite[Tag~07JK]{stacks}, combined with the adjunction between envelopes and the functor forgetting divided powers.

\subsection{Globalisation and coefficients}

The result concerns $\Spec(C)$ only, and passing to schemes over $\Spec(A)$ requires gluing the affine comparison along a cover. Replacing $\Ocris$ by a crystal in quasi-coherent modules leads to the de Rham complex of a module with connection; the present proof treats only the trivial crystal.

\spa

A fully general statement would also have to lift two simplifications made at the outset: the condition $I \cdot C = 0$, in place of the condition that $\gamma$ extends to $C$, and the nilpotence of $p$, in place of $p$-adic completions.

\subsection{Possible contributions to Mathlib}

Several results above do not depend on crystalline cohomology and could be contributed to Mathlib: Lemma~\ref{lem:exterior-quotient}; Propositions~\ref{prop:acyclic-assembly}, \ref{prop:augmented} and~\ref{prop:totmap}, after a comparison with Mathlib's total complexes of bicomplexes; Theorem~\ref{thm:bdj}; Theorem~\ref{thm:contraction}; Theorem~\ref{thm:enough-injectives} and Proposition~\ref{prop:precomposition}; Theorems~\ref{thm:derived-left-adjoint} and~\ref{thm:cech-derived}; and the divided power differentials of Section~\ref{sec:pd-derham}.

%======================================================================
\section*{Acknowledgements}
%======================================================================

This work was supported by the IdEx University of Strasbourg. It was also supported by the Agence Nationale de la Recherche (ANR, French National Research Agency) through the project ANR-25-CE40-7869-01. We thank Anthropic for access to Claude through the Claude Team plan for scientists.

%======================================================================
\appendix
\section{Index of main declarations}\label{app:index}
%======================================================================

Names are given up to their namespace. Most are in \lean{AffineCrystalline} or one of its sub-namespaces; those defined in \lean{PDDifferentials.lean}, \lean{PDDifferentials/}, \lean{SquareZero.lean}, \lean{DeRham.lean}, \lean{DeRham/RiouFull.lean} and \lean{Envelope/PDKaehlerComparison.lean} are in the root namespace. Paths are relative to \lean{AffineCrystalline/}.

{\small
\renewcommand{\arraystretch}{1.2}
\begin{longtable}{@{}>{\raggedright\arraybackslash}p{0.2\linewidth} >{\raggedright\arraybackslash}p{0.44\linewidth} >{\raggedright\arraybackslash}p{0.3\linewidth}@{}}
\toprule
Result & Lean declaration & File \\
\midrule
\endhead
Lemma~\ref{lem:pd-exponent} & \lean{dividedPowers_pow_eq_zero_of_factorial_eq_zero} & \lean{InfinitesimalSite/Envelope.lean} \\
Cor.~\ref{cor:pd-exponent-nilpotent} & \lean{natCast_factorial_pow_eq_zero} & \lean{InfinitesimalSite/Envelope.lean} \\
Prop.~\ref{prop:levels} & \lean{frobeniusExponent_le_iff}, \lean{Level.exists_frobenius_level} & \lean{InfinitesimalSite/Filtration.lean} \\
Prop.~\ref{prop:truncation} & \lean{truncateLift}, \lean{truncate_hom_ext} & \lean{InfinitesimalSite/Filtration.lean} \\
Thm.~\ref{thm:adjunction} & \lean{Level.truncationAdjunction} & \lean{InfinitesimalSite/Filtration.lean} \\
Prop.~\ref{prop:coproducts} & \lean{Level.coprodCofanIsColimit} & \lean{InfinitesimalSite/Filtration.lean} \\
Prop.~\ref{prop:cech-hom} & \lean{Level.cechHomEquiv} & \lean{InfinitesimalSite/Filtration.lean} \\
Prop.~\ref{prop:envelope-unique} & \lean{IsDPEnvelope.comparisonEquiv} & \lean{Envelope/PDKaehlerComparison.lean} \\
Thm.~\ref{thm:envelope-existence} & \lean{DPEnvelope.hasDPEnvelope}, \lean{envelopeExistenceData} & \lean{Envelope/Construction.lean}, \lean{MainTheorems.lean} \\
Prop.~\ref{prop:functoriality} & \lean{EnvelopeChoice.comparison_comp} & \lean{InfinitesimalSite/Envelope.lean} \\
Thm.~\ref{thm:projection} & \lean{projectionKernel_eq_dividedPowersIdeal} & \lean{InfinitesimalSite/Envelope.lean} \\
Lemma~\ref{lem:truncation-envelope} & \lean{EnvelopeChoice.truncate_isInitial}, \lean{EnvelopeFamily.ocrisTruncationIso} & \lean{Ocris/Filtration.lean} \\
Thm.~\ref{thm:adjunction-derived} & \lean{EnvelopeFamily.ocrisDerivedGlobalSectionsIsoLevel} & \lean{Ocris/Filtration.lean} \\
Prop.~\ref{prop:cech-envelope} & \lean{ocrisCech_isEnvelope} & \lean{Ocris/Filtration.lean} \\
Prop.~\ref{prop:pd-kaehler} & \lean{pdKaehlerDerivationEquiv} & \lean{PDDifferentials/Representability.lean} \\
Prop.~\ref{prop:square-zero} & \lean{sqZeroDividedPowers} & \lean{SquareZero.lean} \\
Thm.~\ref{thm:derivation-extension} & \lean{existsUnique_pdDerivation_of_initial} & \lean{Ocris/PDDerivationExtension.lean} \\
Thm.~\ref{thm:envelope-differentials} & \lean{pdKaehler_dpEnvelope} & \lean{PDDifferentials.lean} \\
Lemma~\ref{lem:exterior-quotient} & \lean{exteriorPower_ker_map_mkQ} & \lean{DeRham/ExteriorPowerQuotient.lean} \\
Thm.~\ref{thm:derham} & \lean{DeRhamComplex.d_d} & \lean{DeRham/RiouFull.lean} \\
Thm.~\ref{thm:pd-derham} & \lean{pdQuotientRelationStable}, \lean{pdDeRhamComplex} & \lean{DeRham/PDStabilityFull.lean}, \lean{DeRham.lean} \\
Prop.~\ref{prop:pd-derham-functorial} & \lean{pdDeRhamMap_comp} & \lean{Ocris/DeRham.lean} \\
Prop.~\ref{prop:acyclic-assembly} & \lean{exactAt_totComplex_of_columnsExact} & \lean{Homology/TotM.lean} \\
Prop.~\ref{prop:totmap} & \lean{totMap_quasiIso} & \lean{Homology/TotQuasiIso.lean} \\
Thm.~\ref{thm:bdj} & \lean{RawBi.bhattDeJong_row}, \lean{RawBi.bhattDeJong_column} & \lean{Homology/CosimplicialTot.lean} \\
Thm.~\ref{thm:contraction} & \lean{CosimplicialContraction.exactAt} & \lean{Homology/CosimplicialContraction.lean} \\
Thm.~\ref{thm:enough-injectives} & \lean{hasEnoughInjectivesFunctorCategory} & \lean{Cohomology/FunctorCategoryInjectives.lean} \\
Thm.~\ref{thm:derived-left-adjoint} & \lean{derivedLimitIsoOfAdjunction} & \lean{Ocris/Filtration.lean} \\
Thm.~\ref{thm:cech-zigzag} & \lean{cechZigzag} & \lean{Cohomology/Cech.lean} \\
Thm.~\ref{thm:cech-derived} & \lean{rightDerivedLimIsoCech} & \lean{Cohomology/CechDerived.lean} \\
Thm.~\ref{thm:cech-alexander} & \lean{EnvelopeFamily.crystallineCohomologyIsoCechAlexander} & \lean{Ocris/CechAlexander.lean}, \lean{MainTheorems.lean} \\
Prop.~\ref{prop:presented} & \lean{PresentedEnvelope.eq_top}, \lean{span_dx} & \lean{Ocris/PresentedEnvelope.lean}, \lean{Ocris/PresentedEnvelopeTools.lean} \\
Prop.~\ref{prop:vanishing} & \lean{positiveCechExact} & \lean{Ocris/CechVanishing.lean} \\
Thm.~\ref{thm:one-variable} & \lean{OneVariable.cocycle_succ} & \lean{DeRham/OneVariablePoincare.lean} \\
Prop.~\ref{prop:poincare} & \lean{firstVertexQuasiIso} & \lean{Ocris/CechPoincare.lean} \\
Thm.~\ref{thm:main} & \lean{EnvelopeFamily.crystallineCohomologyIsoDeRham} & \lean{Ocris/CechPoincare.lean}, \lean{MainTheorems.lean} \\
\bottomrule
\end{longtable}
}

%======================================================================
\begin{thebibliography}{99}

\bibitem{stacks}
The Stacks Project Authors, \emph{The Stacks Project}. \url{https://stacks.math.columbia.edu}

\bibitem{bhatt-dejong}
B.~Bhatt and A.~J.~de Jong, \emph{Crystalline cohomology and de Rham cohomology}, preprint, 2011, \href{https://arxiv.org/abs/1110.5001}{arXiv:1110.5001}.

\bibitem{berthelot}
P.~Berthelot, \emph{Cohomologie cristalline des schémas de caractéristique $p > 0$}, Lecture Notes in Mathematics 407, Springer, 1974.

\bibitem{berthelot-ogus}
P.~Berthelot and A.~Ogus, \emph{Notes on crystalline cohomology}, Mathematical Notes 21, Princeton University Press, 1978.

\bibitem{cl-df}
A.~Chambert-Loir and M.~I.~de Frutos-Fernández, \emph{A formalization of divided powers in Lean}, in: 16th International Conference on Interactive Theorem Proving (ITP 2025), LIPIcs, 2025, 4:1--4:17, \href{https://doi.org/10.4230/LIPIcs.ITP.2025.4}{doi:10.4230/LIPIcs.ITP.2025.4}.

\bibitem{cl-df-laws}
A.~Chambert-Loir and M.~I.~de Frutos-Fernández, \emph{Formalizing polynomial laws and the universal divided power algebra}, in: Proceedings of the 15th ACM SIGPLAN International Conference on Certified Programs and Proofs (CPP 2026), ACM, 2026, 75--87, \href{https://doi.org/10.1145/3779031.3779108}{doi:10.1145/3779031.3779108}, \href{https://arxiv.org/abs/2512.05750}{arXiv:2512.05750}.

\bibitem{mathlib}
The mathlib Community, \emph{The Lean mathematical library}, in: Proceedings of the 9th ACM SIGPLAN International Conference on Certified Programs and Proofs (CPP 2020), 2020, 367--381, \href{https://doi.org/10.1145/3372885.3373824}{doi:10.1145/3372885.3373824}. \url{https://github.com/leanprover-community/mathlib4}

\bibitem{roby63}
N.~Roby, \emph{Lois polynômes et lois formelles en théorie des modules}, Ann. Sci. École Norm. Sup. (3) \textbf{80} (1963), 213--348.

\bibitem{roby65}
N.~Roby, \emph{Les algèbres à puissances divisées}, Bull. Sci. Math. (2) \textbf{89} (1965), 75--91.

\bibitem{edged}
M.~D'Addezio, \emph{Edged crystalline cohomology}, preprint, \url{https://daddezio.pages.math.cnrs.fr/edged-crys.pdf}.

\end{thebibliography}

\end{document}
